\documentclass[12pt,a4paper,oneside]{report}
\usepackage{estilo/icmc}
\addbibresource{../../referencias.bib}
\graphicspath{{figuras/}}

% RASCUNHO da dissertacao (nao entra no pre-projeto).
% Compilar dentro desta pasta:
%   portugues: TEXINPUTS=".;../../;" latexmk -pdf -outdir=output -jobname=Dissertacao_MECAI_Felipe_Terrazas main.tex
%   ingles:    TEXINPUTS=".;../../;" latexmk -pdf -outdir=output -jobname=Dissertacao_MECAI_Felipe_Terrazas_EN main.tex
% Tabelas em tabelas/ e figuras em figuras/ sao GERADAS por analise/viabilidade/*.py,
% analise/real/descritiva.py, analise/real/figuras.py e r/descritiva.R:
% nao editar a mao; rodar o script de novo.

\newcommand{\nsitesprev}{18}
\newcommand{\nsemanas}{156}
\newcommand{\fracexposicao}{1}

% Gerado por analise/real/descritiva.py: nao editar a mao.
\newcommand{\descZeros}{\tr{88}{88}}
\newcommand{\descVarMedia}{\tr{1,40}{1.40}}
\newcommand{\descGiniObs}{\tr{0,72}{0.72}}
\newcommand{\descGiniPoisson}{\tr{0,69}{0.69}}
\newcommand{\descGiniHoras}{\tr{0,47}{0.47}}
\newcommand{\descGiniSiteObs}{\tr{0,68}{0.68}}
\newcommand{\descGiniSitePoisson}{\tr{0,65}{0.65}}
\newcommand{\descForaIC}{\tr{8}{8}}
\newcommand{\descPnoventaPdez}{\tr{85}{85}}
\newcommand{\descVelSpotter}{\tr{0,12}{0.12}}
\newcommand{\descVelShuttle}{\tr{0,48}{0.48}}
\newcommand{\descTurnoUm}{\tr{38}{38}}
\newcommand{\descTurnoDois}{\tr{36}{36}}
\newcommand{\descTurnoTres}{\tr{26}{26}}
\newcommand{\descTopCausas}{\tr{41}{41}}
\newcommand{\descSuspeitos}{\tr{24}{24}}
\newcommand{\descComHora}{\tr{97}{97}}
\newcommand{\descHorasExtras}{\tr{10}{10}}
\newcommand{\descNotaAud}{\tr{94}{94}}
\newcommand{\descAcidentesPct}{\tr{87}{87}}

% Gerado por r/descritiva.R: nao editar a mao.
\newcommand{\descRzerosObs}{\tr{88}{88}}
\newcommand{\descRzerosPoisson}{\tr{87}{87}}
\newcommand{\descRzerosNB}{\tr{88}{88}}
\newcommand{\descRzerosRE}{\tr{88}{88}}
\newcommand{\descRveicZeroObs}{\tr{53}{53}}
\newcommand{\descRveicZeroPoisson}{\tr{50}{50}}
\newcommand{\descRveicZeroRE}{\tr{54}{54}}
\newcommand{\descRdpVeiculo}{\tr{0,51}{0.51}}
\newcommand{\descRcorMax}{\tr{0,47}{0.47}}
\newcommand{\descRcorAcidMax}{\tr{0,11}{0.11}}
\newcommand{\descRcorHorasMax}{\tr{0,14}{0.14}}


% \input nao e expansivel dentro de tabular (quebra o \bottomrule); usa o primitivo
\makeatletter\newcommand{\tabinput}[1]{\@@input #1 }\makeatother

\titleformat{\chapter}[display]{\bfseries\color{destaque}\Large}{\tr{Capítulo}{Chapter} \thechapter}{6pt}{\LARGE}
\titlespacing*{\chapter}{0pt}{0pt}{18pt}

\hypersetup{pdftitle={\tr{Dissertação (rascunho)}{Dissertation (draft)}}}

\begin{document}

\begin{titlepage}
\centering
{\large\bfseries \tr{UNIVERSIDADE DE SÃO PAULO}{UNIVERSITY OF SÃO PAULO}\par}
{\large\bfseries \tr{INSTITUTO DE CIÊNCIAS MATEMÁTICAS E DE COMPUTAÇÃO}{INSTITUTE OF MATHEMATICAL AND COMPUTER SCIENCES}\par}
\vspace{4cm}
{\large Felipe Ramires Terrazas\par}
\vspace{2cm}
{\Large\bfseries\color{destaque}\tr%
  {Risco de acidentes em operações de pátio: modelos hierárquicos de contagem com exposição e telemetria veicular}%
  {Accident risk in yard operations: hierarchical count models with exposure and vehicle telematics}\par}
\vspace{1cm}
{\large \tr{Rascunho da dissertação de mestrado (MECAI)}{Master's dissertation draft (MECAI)}\par}
\vspace{0.5cm}
{\normalsize\color{pendencia} \tr{Documento de trabalho, não distribuir}{Working document, do not distribute}\par}
\vfill
{\normalsize São Carlos\par}
{\normalsize \today\par}
\end{titlepage}

\tableofcontents

\chapter{\tr{Introdução}{Introduction}}
\label{cap:introducao}

\section{\tr{Contexto: a operação de pátio}{Context: yard operations}}

\pt{%
Centros de distribuição, fábricas e terminais de carga dependem de uma atividade pouco visível: a movimentação de carretas dentro do pátio. Nos Estados Unidos, onde a carreta costuma ser desacoplada do cavalo mecânico que a trouxe, as carretas esperam no pátio até a hora de carregar ou descarregar. Um caminhão de manobra, chamado de \textit{yard truck}, \textit{yard tractor}, \textit{terminal tractor} ou \textit{spotter}, as leva da vaga para a doca e de volta (\textit{spotting}), e, em operações com mais de uma instalação, entre prédios próximos pelas vias públicas (\textit{shuttling}). É comum que esse serviço seja terceirizado a empresas especializadas, que operam ao mesmo tempo dezenas de pátios de clientes diferentes.
}
\en{%
Distribution centers, plants and freight terminals depend on a barely visible activity: moving trailers inside the yard. In the United States, where the trailer is usually uncoupled from the road tractor that brought it, trailers wait in the yard until it is time to load or unload. A yard truck, also called yard tractor, terminal tractor or spotter, moves them from the parking spot to the dock and back (spotting) and, in operations with more than one facility, between nearby buildings over public roads (shuttling). This service is often outsourced to specialized companies, which run dozens of yards for different clients at the same time.
}

\pt{%
O caminhão de manobra anda pouco e devagar, mas manobra muito: engata e desengata carretas dezenas de vezes por turno, anda de ré até a doca, circula entre carretas estacionadas, pedestres e outros veículos, muitas vezes com pouca visibilidade. O acidente típico não é a colisão em alta velocidade da estrada, mas o choque de baixa velocidade com a carreta, a doca, o portão, uma estrutura fixa ou outro veículo. A maior parte gera só dano material, mas há também lesões e, em casos raros, mortes por atropelamento \parencite{wali2022}. Além do custo do reparo, cada acidente interrompe a operação do cliente.
}
\en{%
The yard truck travels little and slowly, but maneuvers a lot: it couples and uncouples trailers dozens of times per shift, backs up to the dock and moves among parked trailers, pedestrians and other vehicles, often with poor visibility. The typical accident is not a high-speed road crash but a low-speed impact with the trailer, the dock, the gate, a fixed structure or another vehicle. Most cause property damage only, but there are also injuries and, in rare cases, fatal pedestrian strikes \parencite{wali2022}. Besides the repair cost, each accident disrupts the client's operation.
}

\pt{%
O autor trabalha como gestor de inteligência artificial e analytics numa empresa desse tipo, com algumas centenas de caminhões em dezenas de unidades. Nos últimos anos, a empresa passou a ter, ao mesmo tempo, quatro fontes de dados sobre o problema: relatos formais de acidentes, registrados num aplicativo próprio; telemetria embarcada nos caminhões, com viagens, distância, tempo de operação e eventos de condução detectados por sensores; horas trabalhadas por unidade, vindas da folha de pagamento; e ordens de manutenção, que registram os reparos por dano em cada veículo. Essas fontes são analisadas separadamente, e a comparação entre unidades usa contagens brutas de acidentes.
}
\en{%
The author works as head of artificial intelligence and analytics at a company of this kind, with a few hundred trucks at dozens of sites. In recent years the company has come to hold, at the same time, four data sources on the problem: formal accident reports, recorded in an in-house application; on-board telematics in the trucks, with trips, distance, operating time and driving events detected by sensors; hours worked per site, from payroll; and maintenance orders, which record damage repairs on each vehicle. These sources are analyzed separately, and sites are compared by raw accident counts.
}

\section{\tr{O problema: por que contagens brutas enganam}{The problem: why raw counts mislead}}

\pt{%
A contagem bruta de acidentes por unidade é o indicador mais simples e o mais usado, mas responde mal à pergunta que a gestão de segurança faz: onde o risco é maior? Há quatro motivos.
}
\en{%
The raw accident count per site is the simplest and most widely used indicator, but it answers poorly the question safety management asks: where is risk higher? There are four reasons.
}

\pt{%
\begin{enumerate}
  \item \textbf{Exposição.} Unidades com o triplo de caminhões ou de horas operadas tendem a ter mais acidentes sem ser mais arriscadas. A contagem mede sobretudo o tamanho da unidade. O risco é uma taxa: acidentes por unidade de exposição.
  \item \textbf{Acaso.} Acidentes são eventos raros. Numa unidade com um acidente esperado por mês, ter zero ou três num dado mês é perfeitamente compatível com o mesmo risco. Rankings mensais ou trimestrais oscilam por acaso, e quem está no topo num período tende a cair no seguinte sem que nada tenha mudado (regressão à média) \parencite{hauer1997,spiegelhalter2005}.
  \item \textbf{Tipo de operação.} O caminhão que só manobra no pátio, o que faz transferências pela via pública e o de uso local têm riscos e padrões de condução muito diferentes. Comparar unidades com misturas diferentes de operação, ou comparar caminhões pela distância rodada, mistura risco com tipo de serviço.
  \item \textbf{Relato.} Nem todo acidente é relatado, e a propensão a relatar varia com o clima de segurança e a supervisão \parencite{probst2008,probst2010}. Uma unidade com poucos relatos pode ser segura ou pode apenas relatar menos.
\end{enumerate}
}
\en{%
\begin{enumerate}
  \item \textbf{Exposure.} Sites with three times the trucks or operating hours tend to have more accidents without being riskier. The count mostly measures site size. Risk is a rate: accidents per unit of exposure.
  \item \textbf{Chance.} Accidents are rare events. At a site with one expected accident per month, having zero or three in a given month is entirely compatible with the same risk. Monthly or quarterly rankings fluctuate by chance, and whoever is at the top in one period tends to drop in the next without anything having changed (regression to the mean) \parencite{hauer1997,spiegelhalter2005}.
  \item \textbf{Type of operation.} The truck that only maneuvers in the yard, the one that shuttles over public roads and the local-use truck have very different risks and driving patterns. Comparing sites with different operation mixes, or comparing trucks by distance driven, confounds risk with type of service.
  \item \textbf{Reporting.} Not every accident is reported, and the propensity to report varies with safety climate and supervision \parencite{probst2008,probst2010}. A site with few reports may be safe or may simply report less.
\end{enumerate}
}

\pt{%
Os três primeiros problemas têm solução conhecida na estatística: modelos de contagem com a exposição como \textit{offset}, efeitos aleatórios que encolhem as taxas de unidades pequenas em direção à média, e ajuste pelo tipo de operação \parencite{lord2010,gelman2007}. O quarto pede uma segunda fonte que registre o mesmo tipo de evento, para estimar quantos acidentes não chegam aos relatos \parencite{iwgdmf1995}. A telemetria acrescenta uma pergunta própria: se os eventos de condução registrados pelos sensores ajudam a prever acidentes, eles poderiam servir de indicador antecipado de risco, antes que o acidente aconteça.
}
\en{%
The first three problems have known solutions in statistics: count models with exposure as an offset, random effects that shrink the rates of small sites toward the mean, and adjustment for the type of operation \parencite{lord2010,gelman2007}. The fourth calls for a second source that records the same kind of event, to estimate how many accidents never reach the reports \parencite{iwgdmf1995}. Telematics adds a question of its own: if the driving events recorded by the sensors help predict accidents, they could serve as a leading risk indicator, before the accident happens.
}

\section{\tr{Lacuna na literatura}{Gap in the literature}}

\pt{%
A literatura próxima ao tema se divide em grupos que não se tocam. A ciência de dados em logística trata sobretudo de previsão de demanda e otimização, e, quando fala de risco, fala de ruptura da cadeia de suprimentos, não de acidente \parencite{akbari2021,baryannis2019}. Em pátios, o aprendizado de máquina aparece em terminais portuários, com previsão de tempo de permanência de contêineres e de giro de caminhões \parencite{filom2022}, e a pesquisa operacional em centros de distribuição e \textit{cross-docks} estuda o agendamento das docas, não a movimentação no pátio \parencite{gue1999,boysen2010,vanbelle2012}. A segurança de caminhões é estudada na rodovia, onde eventos de condução registrados por telemetria se associam a colisões \parencite{cai2021,driessen2024}, e a segurança em armazéns concentra-se nas empilhadeiras \parencite{collins1999}. O caminhão de manobra fica entre esses grupos, e o que se encontra sobre ele é literatura cinzenta: relatórios de investigação de acidentes, material de fornecedores e patentes \parencite{wali2022}.
}
\en{%
The literature close to the topic splits into groups that do not touch. Data science in logistics deals mostly with demand forecasting and optimization and, when it addresses risk, it means supply chain disruption, not accidents \parencite{akbari2021,baryannis2019}. In yards, machine learning appears in port terminals, forecasting container dwell time and truck turnaround \parencite{filom2022}, and operations research on distribution centers and cross-docks studies dock scheduling, not movement in the yard \parencite{gue1999,boysen2010,vanbelle2012}. Truck safety is studied on the road, where telematics-recorded driving events are associated with crashes \parencite{cai2021,driessen2024}, and warehouse safety focuses on forklifts \parencite{collins1999}. The yard truck falls between these groups, and what one finds about it is grey literature: accident investigation reports, vendor material and patents \parencite{wali2022}.
}

\pt{%
Não foi encontrado estudo revisado por pares que modele o risco de acidentes de caminhões de manobra com exposição medida, nem que teste o valor da telemetria de estrada nesse contexto \preencher{confirmar com busca sistemática em Scopus e Web of Science, com as strings de \texttt{notas/06-mapa-da-literatura.md}, seção 6; as buscas na web de 2026-11-08 e 2026-11-12 não acharam nenhum}. Essa é a lacuna que o trabalho ocupa.
}
\en{%
No peer-reviewed study was found that models the accident risk of yard trucks with measured exposure, or that tests the value of road telematics in this setting \preencher{confirm with a systematic search in Scopus and Web of Science}. This is the gap the work fills.
}

\section{\tr{Contribuições}{Contributions}}

\pt{%
\begin{itemize}
  \item \textbf{Medição do risco.} Um arcabouço para medir e comparar o risco de acidentes entre unidades e veículos de pátio, com exposição medida pela telemetria, ajuste pelo tipo de operação e incerteza explícita, e uma demonstração de que a escolha da medida de exposição muda as conclusões.
  \item \textbf{Telemetria.} Um teste, com testes de equivalência para o caso nulo, de quanto os eventos de condução de uma telemetria desenhada para a estrada acrescentam à previsão de acidentes no pátio, e uma medida da cobertura da detecção automática de colisão nesse contexto.
  \item \textbf{Subnotificação.} Uma estimativa, com limites explícitos, de quanto os relatos subestimam os acidentes e de como isso varia entre unidades, usando as ordens de manutenção como segunda fonte.
  \item \textbf{Método.} Um estudo de simulação calibrado em dados públicos que mostra quantos acidentes e que estrutura de modelo são necessários para essas perguntas em frotas de algumas centenas de veículos.
  \item \textbf{Aplicação.} Um indicador de risco pronto para uso pela gestão de segurança e, como extensão, uma previsão probabilística das horas de operação para o planejamento de escala e supervisão.
\end{itemize}
}
\en{%
\begin{itemize}
  \item \textbf{Risk measurement.} A framework to measure and compare accident risk across yard sites and vehicles, with exposure measured by telematics, adjustment for type of operation and explicit uncertainty, and a demonstration that the choice of exposure measure changes the conclusions.
  \item \textbf{Telematics.} A test, with equivalence tests for the null case, of how much the driving events of road-oriented telematics add to accident prediction in the yard, and a measure of the coverage of automatic crash detection in this setting.
  \item \textbf{Underreporting.} An estimate, with explicit bounds, of how much reports underestimate accidents and how this varies across sites, using maintenance orders as a second source.
  \item \textbf{Method.} A simulation study calibrated on public data showing how many accidents and which model structure these questions require in fleets of a few hundred vehicles.
  \item \textbf{Application.} A risk indicator ready for use by safety management and, as an extension, a probabilistic forecast of operating hours for staffing and supervision planning.
\end{itemize}
}

\section{\tr{Objetivos}{Objectives}}
\label{sec:objetivos}

\subsection{\tr{Objetivo geral}{General objective}}

\pt{%
Desenvolver e avaliar modelos estatísticos para estimar o risco de acidentes em operações de pátio, ajustado pela exposição e com incerteza quantificada, por veículo e por unidade operacional, combinando relatos de acidentes, telemetria veicular, horas trabalhadas e ordens de manutenção.
}
\en{%
To develop and evaluate statistical models that estimate accident risk in yard operations, adjusted for exposure and with quantified uncertainty, per vehicle and per operating site, by combining accident reports, vehicle telematics, hours worked and maintenance orders.
}

\subsection{\tr{Objetivos específicos}{Specific objectives}}

\begin{enumerate}[label=(\alph*)]
  \item \tr%
    {Integrar as fontes num painel veículo $\times$ mês, com regras de ligação documentadas e qualidade da ligação medida em amostra revisada manualmente.}%
    {Integrate the sources into a vehicle $\times$ month panel, with documented linkage rules and linkage quality measured on a manually reviewed sample.}
  \item \tr%
    {Comparar modelos de contagem (Poisson, binomial negativa, com e sem inflação de zeros, com e sem efeitos aleatórios de veículo e unidade) quanto ao ajuste e ao desempenho preditivo fora da amostra.}%
    {Compare count models (Poisson, negative binomial, with and without zero inflation, with and without vehicle and site random effects) in terms of fit and out-of-sample predictive performance.}
  \item \tr%
    {Medir quanto as taxas de eventos de condução registradas pela telemetria acrescentam à previsão de acidentes, além da exposição, do tipo de uso do veículo e da unidade, e, se nada acrescentarem, verificar se o efeito pode ser declarado desprezível.}%
    {Measure how much the driving event rates recorded by telematics add to accident prediction beyond exposure, vehicle use type and site, and, if they add nothing, check whether the effect can be declared negligible.}
  \item \tr%
    {Estimar a subnotificação comparando os relatos formais com as ordens de manutenção abertas por dano, e medir a cobertura da detecção automática de colisão da telemetria como fonte de acidentes de pátio.}%
    {Estimate underreporting by comparing formal reports with maintenance orders opened for damage, and measure the coverage of telematics automatic crash detection as a source of yard accidents.}
  \item \tr%
    {Produzir indicadores de risco por unidade ajustados por exposição, com intervalos de incerteza e grupos de risco, e compará-los aos indicadores por contagem bruta quanto à capacidade de prever o período seguinte e à confusão com o tamanho da unidade.}%
    {Produce exposure-adjusted site risk indicators, with uncertainty intervals and risk groups, and compare them with raw-count indicators in their ability to predict the next period and in their confounding with site size.}
  \item \textit{\tr{(Opcional, extensão)}{(Optional, extension)}} \tr%
    {Prever de forma probabilística a exposição (horas de operação) de cada unidade nos meses seguintes, como insumo para o planejamento de escala e de supervisão, e combinar essa previsão ao modelo de risco para projetar acidentes esperados em níveis agregados (região e trimestre).}%
    {Forecast each site's exposure (operating hours) for the following months probabilistically, as input for staffing and supervision planning, and combine this forecast with the risk model to project expected accidents at aggregate levels (region and quarter).}
\end{enumerate}

\subsection{\tr{Hipóteses}{Hypotheses}}
\label{sec:hipoteses}

\pt{%
Cada hipótese de pesquisa é testada contra uma hipótese nula de igualdade ($H_0$), de modo que rejeitar $H_0$ permite afirmar a diferença. Como não rejeitar $H_0$ não demonstra igualdade, as afirmações de ausência de efeito (H1c e H2b) usam testes de equivalência \parencite{schuirmann1987,lakens2017}: nesses casos, $H_0$ é que o efeito está fora de uma margem fixada antes da análise final, e rejeitá-la permite afirmar que o efeito é desprezível. Os procedimentos de teste estão na Seção~\ref{sec:testes}.
}
\en{%
Each research hypothesis is tested against a null hypothesis of equality ($H_0$), so that rejecting $H_0$ supports the claim of a difference. Since failing to reject $H_0$ does not demonstrate equality, claims of no effect (H1c and H2b) use equivalence tests \parencite{schuirmann1987,lakens2017}: in those cases, $H_0$ is that the effect lies outside a margin fixed before the final analysis, and rejecting it supports the claim that the effect is negligible. The test procedures are in Section~\ref{sec:testes}.
}

\begin{description}[leftmargin=1.2cm,style=nextline]
  \item[H1] \tr%
    {\textit{Comparação entre unidades.} (a)~$H_0$: o erro de previsão da taxa do semestre seguinte é igual com a taxa estimada pelo modelo hierárquico e com a taxa bruta; alternativa: é menor com a estimada. (b)~$H_0$: a taxa de acidentes é igual em todas as unidades, dada a exposição; alternativa: difere entre unidades. (c)~Equivalência: $H_0$: a correlação de postos entre a taxa estimada e a exposição da unidade está fora de $(-0{,}3;\ 0{,}3)$; alternativa: está dentro, ou seja, a taxa estimada não se confunde com o tamanho da unidade.}%
    {\textit{Site comparison.} (a)~$H_0$: the error in predicting the next semester's rate is the same with the hierarchical model's rate and with the raw rate; alternative: it is smaller with the model's rate. (b)~$H_0$: the accident rate is the same in all sites, given exposure; alternative: it differs across sites. (c)~Equivalence: $H_0$: the rank correlation between the estimated rate and site exposure lies outside $(-0.3,\ 0.3)$; alternative: it lies inside, that is, the estimated rate is not confounded with site size.}
  \item[H2] \tr%
    {\textit{Telemetria.} (a)~$H_0$: para cada tipo de evento de condução, a razão de taxas de acidente por desvio padrão da taxa de eventos é igual a 1, dados exposição, tipo de uso e veículo; e o escore logarítmico fora da amostra é igual com e sem telemetria; alternativa: diferente de 1, e escore maior com telemetria. (b)~Equivalência: $H_0$: a razão de taxas está fora de $[0{,}8;\ 1{,}25]$; alternativa: está dentro, ou seja, o efeito é desprezível.}%
    {\textit{Telematics.} (a)~$H_0$: for each type of driving event, the accident rate ratio per standard deviation of the event rate equals 1, given exposure, use type and vehicle; and the out-of-sample log score is the same with and without telematics; alternative: different from 1, and a higher score with telematics. (b)~Equivalence: $H_0$: the rate ratio lies outside $[0.8,\ 1.25]$; alternative: it lies inside, that is, the effect is negligible.}
  \item[H3] \tr%
    {\textit{Subnotificação.} (a)~$H_0$: a probabilidade de que um dano registrado pela manutenção tenha relato formal correspondente é igual a 1; alternativa: é menor que 1. (b)~$H_0$: essa probabilidade é igual em todas as unidades; alternativa: varia entre unidades, de modo que as contagens de relatos subestimam o risco de forma desigual.}%
    {\textit{Underreporting.} (a)~$H_0$: the probability that damage recorded by maintenance has a matching formal report equals 1; alternative: it is below 1. (b)~$H_0$: this probability is the same in all sites; alternative: it varies across sites, so that report counts underestimate risk unevenly.}
\end{description}

\subsection{\tr{O que mudou desde o pré-projeto}{What changed since the pre-proposal}}

\pt{%
Os estudos preliminares do Capítulo~\ref{cap:viabilidade} já alteraram o desenho original em cinco pontos: (i) o desfecho principal passou a ser o conjunto de todas as colisões relatadas, com ou sem lesão, porque uma definição estreita derruba o poder de H2; (ii) o efeito aleatório de veículo deixou de ser opcional, porque sem ele a inferência sobre telemetria tem falsos positivos acima do nominal; (iii) o tipo de uso do veículo entrou como ajuste obrigatório, e o tempo de viagem substituiu a distância como exposição principal; (iv) a segunda fonte para a subnotificação passou a ser a manutenção, porque a detecção automática de colisão quase não capta acidentes de pátio; (v) a extensão (f) passou a ter como uso principal o planejamento de escala, com projeção de acidentes só em níveis agregados.
}
\en{%
The preliminary studies in Chapter~\ref{cap:viabilidade} have already changed the original design in five ways: (i) the main outcome became all reported collisions, with or without injury, because a narrow definition kills the power of H2; (ii) the vehicle random effect is no longer optional, because without it inference on telematics has false positives above the nominal level; (iii) vehicle use type became a mandatory adjustment, and trip time replaced distance as the main exposure; (iv) the second source for underreporting became maintenance, because automatic crash detection captures almost no yard accidents; (v) extension (f) now has staffing planning as its main use, with accident projection only at aggregate levels.
}

\section{\tr{Organização do texto}{Outline}}

\pt{%
O Capítulo~\ref{cap:fundamentacao} revisa a literatura em três blocos: o contexto (ciência de dados em logística, pátios e terminais, segurança de caminhões e segurança ocupacional), os métodos estatísticos (modelos de contagem com exposição, modelos hierárquicos e Bayes empírico, captura e recaptura, avaliação de previsões e testes de equivalência) e a lacuna que motiva o trabalho. O Capítulo~\ref{cap:dados} descreve as fontes de dados da operação, as definições de desfecho, exposição e tipo de uso, os problemas de qualidade, e o conjunto público usado nos estudos preliminares. A análise segue três níveis. O Capítulo~\ref{cap:descritiva} é descritivo: mostra a distribuição de cada métrica usada nos modelos (exposição, desfecho, concentração, tempo, causas, telemetria e fontes de subnotificação) e o que ela implica para a modelagem. O Capítulo~\ref{cap:metodos} detalha a construção do painel, o modelo binomial negativo hierárquico e as alternativas comparadas, a validação fora da amostra, os procedimentos de teste de cada hipótese, a estimação da subnotificação, o plano alternativo e a extensão de previsão. O Capítulo~\ref{cap:viabilidade} apresenta os estudos preliminares: uma simulação calibrada em dados públicos, que mede o poder e os falsos positivos dos testes, e uma exploração agregada dos dados da operação, que orientou as revisões listadas acima. O Capítulo~\ref{cap:preditiva} é preditivo: ajusta os modelos com os dados completos e testa as hipóteses. O Capítulo~\ref{cap:prescritiva} é prescritivo: transforma as estimativas em regras de decisão (priorização de unidades, escala de motoristas por turno e alocação de auditorias) e as avalia retrospectivamente. O Capítulo~\ref{cap:conclusao} discute contribuições, limitações e trabalhos futuros.
}
\en{%
Chapter~\ref{cap:fundamentacao} reviews the literature in three blocks: the context (data science in logistics, yards and terminals, truck safety and occupational safety), the statistical methods (count models with exposure, hierarchical models and empirical Bayes, capture-recapture, forecast evaluation and equivalence tests) and the gap that motivates the work. Chapter~\ref{cap:dados} describes the operational data sources, the definitions of outcome, exposure and use type, the quality issues, and the public dataset used in the preliminary studies. The analysis has three levels. Chapter~\ref{cap:descritiva} is descriptive: it shows the distribution of each metric used in the models (exposure, outcome, concentration, time, causes, telematics and underreporting sources) and what it implies for modeling. Chapter~\ref{cap:metodos} details the panel construction, the hierarchical negative binomial model and the alternatives compared, out-of-sample validation, the test procedures for each hypothesis, underreporting estimation, the fallback plan and the forecasting extension. Chapter~\ref{cap:viabilidade} presents the preliminary studies: a simulation calibrated on public data, which measures the power and false positives of the tests, and an aggregated exploration of the operational data, which guided the revisions listed above. Chapter~\ref{cap:preditiva} is predictive: it fits the models with the full data and tests the hypotheses. Chapter~\ref{cap:prescritiva} is prescriptive: it turns the estimates into decision rules (site prioritization, driver staffing by shift and audit allocation) and evaluates them retrospectively. Chapter~\ref{cap:conclusao} discusses contributions, limitations and future work.
}

\chapter{\tr{Fundamentação teórica e revisão}{Background and literature review}}
\label{cap:fundamentacao}

\pt{%
O capítulo tem três partes. As Seções~\ref{sec:fund-logistica} a \ref{sec:fund-ocupacional} situam o problema na literatura de logística e de segurança. As Seções~\ref{sec:fund-contagem} a \ref{sec:fund-testes} apresentam os métodos estatísticos usados no trabalho. A Seção~\ref{sec:fund-lacuna} resume a lacuna.
}
\en{%
The chapter has three parts. Sections~\ref{sec:fund-logistica} to \ref{sec:fund-ocupacional} place the problem in the logistics and safety literature. Sections~\ref{sec:fund-contagem} to \ref{sec:fund-testes} present the statistical methods used in the work. Section~\ref{sec:fund-lacuna} summarizes the gap.
}

\section{\tr{Ciência de dados em logística}{Data science in logistics}}
\label{sec:fund-logistica}

\pt{%
Revisões da aplicação de aprendizado de máquina e de \textit{analytics} em logística e cadeias de suprimentos classificam os trabalhos em descritivos, preditivos e prescritivos, e em estratégicos e operacionais \parencite{wang2016,nguyen2018}. A maior parte trata de previsão de demanda, roteirização e estoques; aplicações de detecção e de segurança aparecem pouco \parencite{akbari2021}. Quando essa literatura fala de risco, o sentido é o de ruptura da cadeia (atraso de fornecedor, falta de material, eventos externos), não o de acidente com pessoas ou equipamentos \parencite{baryannis2019}. Na classificação dessas revisões, este trabalho é preditivo e operacional, e o risco estudado é o de acidente.
}
\en{%
Reviews of machine learning and analytics in logistics and supply chains classify work as descriptive, predictive or prescriptive, and as strategic or operational \parencite{wang2016,nguyen2018}. Most of it addresses demand forecasting, routing and inventory; detection and safety applications are rare \parencite{akbari2021}. When this literature speaks of risk, it means supply chain disruption (supplier delay, material shortage, external events), not accidents involving people or equipment \parencite{baryannis2019}. In the terms of these reviews, this work is predictive and operational, and the risk studied is accident risk.
}

\section{\tr{Pátios, docas e terminais}{Yards, docks and terminals}}
\label{sec:fund-patios}

\pt{%
O pátio aparece na literatura sobretudo pelo lado da eficiência. Em terminais portuários, uma revisão sistemática encontrou dezenas de aplicações de aprendizado de máquina, com muitos dados subutilizados e pouco uso em tempo real \parencite{filom2022}; exemplos são a previsão do tempo de permanência de contêineres \parencite{chhetri2026} e os sistemas de agendamento de caminhões na portaria \parencite{huynh2016}. Em centros de distribuição, há estudos de campo sobre a previsão do horário de chegada de caminhões \parencite{vanderspoel2017} e do tempo de permanência em instalações de um operador logístico \parencite{benjatanont2020}. A pesquisa operacional em terminais de carga fracionada e \textit{cross-docks} estuda a alocação de carretas às docas e o layout \parencite{gue1999,boysen2010,vanbelle2012}.
}
\en{%
The yard appears in the literature mostly from the efficiency side. In port terminals, a systematic review found dozens of machine learning applications, with much underused data and little real-time use \parencite{filom2022}; examples are container dwell time prediction \parencite{chhetri2026} and truck appointment systems at the gate \parencite{huynh2016}. In distribution centers, there are field studies on forecasting truck arrival times \parencite{vanderspoel2017} and dwell time at the facilities of a logistics provider \parencite{benjatanont2020}. Operations research on less-than-truckload terminals and cross-docks studies the assignment of trailers to docks and layout \parencite{gue1999,boysen2010,vanbelle2012}.
}

\pt{%
O elo entre pátio e segurança aparece de forma indireta. Um relatório do órgão de auditoria do Departamento de Transportes dos EUA associou o tempo de espera dos motoristas nas instalações (\textit{detention}) a um risco maior de acidentes na estrada e apontou que os dados nacionais sobre esse tempo são insuficientes \parencite{dotoig2018}. Sobre o acidente dentro do pátio, com o caminhão de manobra, o que se encontra são relatórios de investigação de casos fatais, como o de um motorista atropelado por um caminhão de pátio à noite, com pouca iluminação \parencite{wali2022}.
}
\en{%
The link between yard and safety appears indirectly. A report by the US Department of Transportation's audit office associated driver waiting time at facilities (detention) with higher crash risk on the road and noted that national data on this time are insufficient \parencite{dotoig2018}. On accidents inside the yard, involving the yard truck, what one finds are investigation reports of fatal cases, such as that of a driver struck by a yard truck at night, in poor lighting \parencite{wali2022}.
}

\section{\tr{Telemetria e risco em caminhões}{Telematics and truck risk}}
\label{sec:fund-telemetria}

\pt{%
Sistemas de telemetria registram eventos cinemáticos (acelerações, frenagens e curvas acima de limiares) e eventos de comportamento (velocidade acima do limite, parada incompleta, uso de cinto ou de celular). Estudos de condução naturalística, com veículos instrumentados, associam a frequência desses eventos ao risco de colisão de cada motorista \parencite{guo2013} e quantificam o peso de fatores como distração e fadiga \parencite{dingus2016}. Para caminhões, a evidência mais próxima deste trabalho vem de uma grande transportadora dos EUA, em que eventos críticos de segurança registrados pela telemetria, por unidade de distância, se associaram a acidentes num modelo binomial negativo com a distância como exposição \parencite{cai2021}; o mesmo grupo modelou a ocorrência desses eventos ao longo do tempo com processos pontuais hierárquicos \parencite{cai2022}. Na Holanda, frenagens bruscas por hora de condução previram sinistros com dano material, o desfecho mais comum no pátio \parencite{driessen2024}. Revisões de análise de dados em segurança viária situam esses trabalhos entre os modelos de previsão de risco \parencite{mehdizadeh2020}, e o debate sobre os escores de segurança de transportadoras nos EUA mostra o problema de indicadores que não ajustam bem pela exposição \parencite{nas2017}.
}
\en{%
Telematics systems record kinematic events (accelerations, braking and turns above thresholds) and behavioral events (speeding, rolling stops, seat belt or phone use). Naturalistic driving studies, with instrumented vehicles, associate the frequency of these events with each driver's crash risk \parencite{guo2013} and quantify the weight of factors such as distraction and fatigue \parencite{dingus2016}. For trucks, the evidence closest to this work comes from a large US carrier, where telematics-recorded safety-critical events per unit distance were associated with crashes in a negative binomial model with distance as exposure \parencite{cai2021}; the same group modeled the occurrence of these events over time with hierarchical point processes \parencite{cai2022}. In the Netherlands, harsh braking per driving hour predicted property damage incidents, the most common outcome in the yard \parencite{driessen2024}. Reviews of data analysis in road safety place this work among risk prediction models \parencite{mehdizadeh2020}, and the debate on US carrier safety scores shows the problem of indicators that do not adjust well for exposure \parencite{nas2017}.
}

\pt{%
Esses resultados vêm da estrada. No pátio, as velocidades são baixas, a ré e o engate dominam, e os limiares dos sensores foram calibrados para outro tipo de condução. Não é óbvio que os mesmos eventos carreguem sinal, nem que a detecção automática de colisão, desenhada para impactos de estrada, capte o choque de baixa velocidade típico do pátio. Há também uma distinção de objetivo: um coeficiente significativo num modelo explicativo não garante ganho de previsão fora da amostra, e vice-versa \parencite{shmueli2010}. Por isso H2 é testada pelas duas vias: o coeficiente e o escore preditivo.
}
\en{%
These results come from the road. In the yard, speeds are low, reversing and coupling dominate, and sensor thresholds were calibrated for another kind of driving. It is not obvious that the same events carry signal, nor that automatic crash detection, designed for road impacts, captures the low-speed impact typical of the yard. There is also a distinction of purpose: a significant coefficient in an explanatory model does not guarantee an out-of-sample predictive gain, and vice versa \parencite{shmueli2010}. That is why H2 is tested both ways: the coefficient and the predictive score.
}

\section{\tr{Segurança ocupacional e fatores de risco}{Occupational safety and risk factors}}
\label{sec:fund-ocupacional}

\pt{%
Uma revisão de aprendizado de máquina na análise de acidentes de trabalho organiza a área em previsão de desfecho, descoberta de padrões, avaliação de risco e gravidade \parencite{sarkar2020}. Em armazéns e plantas industriais, o estudo epidemiológico clássico de veículos industriais motorizados trata de empilhadeiras \parencite{collins1999}. Fatores de organização do trabalho conhecidos por aumentar o risco de lesão são candidatos naturais a covariáveis por unidade e mês: horas extras e jornadas longas \parencite{dembe2005}, pouco tempo de casa \parencite{breslin2006} e trabalho temporário \parencite{smith2010}. Os dois primeiros podem ser calculados a partir da folha de pagamento, e entram como análise secundária (Seção~\ref{sec:met-sensibilidade}).
}
\en{%
A review of machine learning in occupational accident analysis organizes the field into outcome prediction, pattern discovery, risk assessment and severity \parencite{sarkar2020}. In warehouses and plants, the classic epidemiological study of powered industrial vehicles deals with forklifts \parencite{collins1999}. Work organization factors known to raise injury risk are natural candidate covariates per site and month: overtime and long hours \parencite{dembe2005}, short job tenure \parencite{breslin2006} and temporary work \parencite{smith2010}. The first two can be computed from payroll and enter as a secondary analysis (Section~\ref{sec:met-sensibilidade}).
}

\section{\tr{Modelos de contagem com exposição}{Count models with exposure}}
\label{sec:fund-contagem}

\pt{%
Seja $Y_{it}$ o número de acidentes do veículo $i$ no mês $t$ e $E_{it}$ sua exposição (horas de operação ou distância). O modelo de Poisson com \textit{offset} supõe
}
\en{%
Let $Y_{it}$ be the number of accidents of vehicle $i$ in month $t$ and $E_{it}$ its exposure (operating hours or distance). The Poisson model with an offset assumes
}
\begin{equation}
\label{eq:poisson}
Y_{it} \sim \text{Poisson}(\mu_{it}), \qquad \log \mu_{it} = \log E_{it} + \mathbf{x}_{it}^\top \bm{\beta},
\end{equation}
\pt{%
de modo que $\lambda_{it} = \mu_{it}/E_{it} = \exp(\mathbf{x}_{it}^\top \bm{\beta})$ é uma taxa por unidade de exposição e $\exp(\beta_j)$ é uma razão de taxas \parencite{cameron2013}. O \textit{offset} impõe que, mantido o resto, dobrar a exposição dobra o número esperado de acidentes. Essa hipótese pode ser verificada incluindo $\log E_{it}$ como covariável com coeficiente livre: um coeficiente perto de 1 a sustenta, e um coeficiente menor indica que a medida de exposição cresce mais rápido que o risco.
}
\en{%
so that $\lambda_{it} = \mu_{it}/E_{it} = \exp(\mathbf{x}_{it}^\top \bm{\beta})$ is a rate per unit of exposure and $\exp(\beta_j)$ is a rate ratio \parencite{cameron2013}. The offset imposes that, all else equal, doubling exposure doubles the expected number of accidents. This assumption can be checked by including $\log E_{it}$ as a covariate with a free coefficient: a coefficient near 1 supports it, and a smaller one indicates that the exposure measure grows faster than risk.
}

\pt{%
No Poisson, a variância é igual à média. Dados de acidentes costumam ter variância maior (superdispersão), porque veículos e períodos diferem em risco por fatores não medidos. A binomial negativa resulta de supor que a taxa de Poisson tem uma variação gama multiplicativa, e na parametrização usual (NB2) tem
}
\en{%
In the Poisson model, the variance equals the mean. Accident data usually show a larger variance (overdispersion), because vehicles and periods differ in risk due to unmeasured factors. The negative binomial results from assuming a multiplicative gamma variation in the Poisson rate, and in the usual parameterization (NB2) has
}
\begin{equation}
\label{eq:nb2}
\mathbb{E}(Y_{it}) = \mu_{it}, \qquad \text{Var}(Y_{it}) = \mu_{it} + \mu_{it}^2/\phi,
\end{equation}
\pt{%
em que $\phi > 0$ controla a dispersão e o Poisson é o limite $\phi \to \infty$ \parencite{cameron2013}. A literatura de frequência de acidentes adota a binomial negativa como ponto de partida e discute extensões com efeitos aleatórios, parâmetros aleatórios e estruturas espaciais e temporais \parencite{lord2010,mannering2014}.
}
\en{%
where $\phi > 0$ controls dispersion and the Poisson is the limit $\phi \to \infty$ \parencite{cameron2013}. The crash frequency literature adopts the negative binomial as the starting point and discusses extensions with random effects, random parameters and spatial and temporal structures \parencite{lord2010,mannering2014}.
}

\pt{%
Quando há mais zeros do que a binomial negativa prevê, uma opção é o modelo com inflação de zeros, que mistura uma massa em zero com a distribuição de contagem \parencite{lambert1992}. Em segurança viária, esse modelo foi criticado por supor um estado ``seguro'', com risco exatamente nulo, que não tem interpretação física: o excesso de zeros costuma vir de exposição pequena, períodos curtos ou heterogeneidade não modelada, e se resolve melhor com a exposição correta e efeitos aleatórios \parencite{lord2005}. Por isso a inflação de zeros entra aqui só como alternativa a ser testada, e não como modelo principal.
}
\en{%
When there are more zeros than the negative binomial predicts, one option is the zero-inflated model, which mixes a point mass at zero with the count distribution \parencite{lambert1992}. In road safety this model has been criticized for assuming a ``safe'' state with exactly zero risk, which has no physical interpretation: excess zeros usually come from small exposure, short periods or unmodeled heterogeneity, and are better handled with the right exposure and random effects \parencite{lord2005}. That is why zero inflation enters here only as an alternative to be tested, not as the main model.
}

\section{\tr{Modelos hierárquicos, encolhimento e Bayes empírico}{Hierarchical models and empirical Bayes}}
\label{sec:fund-hierarquicos}

\pt{%
Os dados têm estrutura aninhada: meses dentro de veículos e veículos dentro de unidades. Efeitos aleatórios de veículo ($v_i$) e de unidade ($u_s$) acrescentados ao preditor linear capturam a correlação entre observações do mesmo grupo e permitem estimar a variação de risco entre grupos ($\sigma_v^2$, $\sigma_u^2$) \parencite{gelman2007}. Ignorar essa correlação trata meses do mesmo veículo como independentes e subestima os erros-padrão das covariáveis que variam entre veículos; a alternativa sem efeitos aleatórios são erros-padrão robustos por agrupamento \parencite{cameron2015}.
}
\en{%
The data have a nested structure: months within vehicles and vehicles within sites. Vehicle ($v_i$) and site ($u_s$) random effects added to the linear predictor capture the correlation between observations in the same group and allow estimating the risk variation across groups ($\sigma_v^2$, $\sigma_u^2$) \parencite{gelman2007}. Ignoring this correlation treats months of the same vehicle as independent and underestimates the standard errors of covariates that vary across vehicles; the alternative without random effects is cluster-robust standard errors \parencite{cameron2015}.
}

\pt{%
O efeito aleatório produz encolhimento: a estimativa da taxa de cada grupo fica entre a taxa observada no grupo e a média geral, mais perto da observada quanto mais informação o grupo tem. A ideia é a mesma do estimador de Stein e do Bayes empírico \parencite{efron1975}. No caso mais simples, com taxas de unidade $\lambda_s$ vindas de uma gama com média $m$ e variância $v$, e $Y_s \mid \lambda_s \sim \text{Poisson}(E_s \lambda_s)$, a média a posteriori é
}
\en{%
The random effect produces shrinkage: each group's rate estimate lies between the group's observed rate and the overall mean, closer to the observed one the more information the group has. The idea is the same as Stein's estimator and empirical Bayes \parencite{efron1975}. In the simplest case, with site rates $\lambda_s$ drawn from a gamma with mean $m$ and variance $v$, and $Y_s \mid \lambda_s \sim \text{Poisson}(E_s \lambda_s)$, the posterior mean is
}
\begin{equation}
\label{eq:eb}
\hat{\lambda}_s^{\text{BE}} = \frac{Y_s + m^2/v}{E_s + m/v} = w_s \, \frac{Y_s}{E_s} + (1 - w_s)\, m, \qquad w_s = \frac{E_s}{E_s + m/v},
\end{equation}
\pt{%
e $m$ e $v$ são estimados com os dados de todas as unidades, daí o nome Bayes empírico. Unidades com muita exposição têm peso $w_s$ perto de 1 e ficam com a própria taxa; unidades pequenas são puxadas para a média. O método é usado há décadas no mapeamento de doenças \parencite{clayton1987} e na segurança viária, onde corrige a regressão à média na escolha de locais críticos e na avaliação antes e depois de intervenções \parencite{hauer1997,hauer2002}.
}
\en{%
and $m$ and $v$ are estimated from the data of all sites, hence the name empirical Bayes. Sites with large exposure have weight $w_s$ near 1 and keep their own rate; small sites are pulled toward the mean. The method has been used for decades in disease mapping \parencite{clayton1987} and in road safety, where it corrects regression to the mean when selecting hazardous locations and evaluating interventions before and after \parencite{hauer1997,hauer2002}.
}

\pt{%
Comparar instituições por rankings tem limites conhecidos: a incerteza sobre a posição de cada unidade costuma ser tão grande que só os extremos se distinguem com segurança, e os intervalos das posições devem acompanhar qualquer tabela \parencite{goldstein1996}. O gráfico de funil é uma alternativa ao ranking: a taxa de cada unidade é plotada contra sua exposição, com limites de controle que se estreitam à medida que a exposição cresce, e só as unidades fora dos limites são destacadas \parencite{spiegelhalter2005}.
}
\en{%
Comparing institutions by rankings has known limits: the uncertainty about each unit's position is usually so large that only the extremes can be told apart with confidence, and position intervals should accompany any table \parencite{goldstein1996}. The funnel plot is an alternative to ranking: each site's rate is plotted against its exposure, with control limits that narrow as exposure grows, and only sites outside the limits are flagged \parencite{spiegelhalter2005}.
}

\pt{%
Modelos lineares generalizados mistos com binomial negativa podem ser ajustados por máxima verossimilhança com a integral dos efeitos aleatórios aproximada pelo método de Laplace, como no pacote glmmTMB \parencite{brooks2017}, construído sobre diferenciação automática \parencite{kristensen2016}, ou de forma bayesiana, com amostragem de Monte Carlo \parencite{burkner2017}.
}
\en{%
Negative binomial generalized linear mixed models can be fitted by maximum likelihood with the random effects integral approximated by the Laplace method, as in the glmmTMB package \parencite{brooks2017}, built on automatic differentiation \parencite{kristensen2016}, or in a Bayesian way, with Monte Carlo sampling \parencite{burkner2017}.
}

\section{\tr{Subnotificação e captura e recaptura}{Underreporting and capture-recapture}}
\label{sec:fund-subnotificacao}

\pt{%
A subnotificação de lesões de trabalho é grande e bem documentada. Nos EUA, estimativas indicam que o levantamento oficial deixa de contar uma parte considerável das lesões não fatais \parencite{leigh2004}, e estudos de captura e recaptura que cruzam o levantamento com outros registros chegaram a conclusões semelhantes, com grande variação entre estados \parencite{rosenman2006,boden2008}. Dentro das empresas, a subnotificação é maior onde o clima de segurança é pior e a supervisão cobra menos as práticas seguras \parencite{probst2008,probst2010}. Para este trabalho, a consequência é que a taxa de relato pode variar entre unidades e distorcer qualquer comparação feita só com relatos (H3b).
}
\en{%
Underreporting of work injuries is large and well documented. In the US, estimates indicate that the official survey misses a considerable share of nonfatal injuries \parencite{leigh2004}, and capture-recapture studies that cross-match the survey with other records reached similar conclusions, with large variation across states \parencite{rosenman2006,boden2008}. Within companies, underreporting is higher where safety climate is worse and supervisors enforce safe practices less \parencite{probst2008,probst2010}. For this work, the consequence is that the reporting rate may vary across sites and distort any comparison based only on reports (H3b).
}

\pt{%
Quando o mesmo tipo de evento é registrado por duas fontes imperfeitas, a sobreposição entre elas informa quantos eventos nenhuma das duas captou \parencite{iwgdmf1995,hook1995}. Com $n_1$ eventos na fonte 1, $n_2$ na fonte 2 e $m$ em ambas, o estimador de Lincoln e Petersen do total é
}
\en{%
When the same type of event is recorded by two imperfect sources, the overlap between them informs how many events neither captured \parencite{iwgdmf1995,hook1995}. With $n_1$ events in source 1, $n_2$ in source 2 and $m$ in both, the Lincoln-Petersen estimator of the total is
}
\begin{equation}
\label{eq:lp}
\hat{N} = \frac{n_1 \, n_2}{m},
\end{equation}
\pt{%
e $m/n_2$ estima a probabilidade de captura da fonte 1. O estimador supõe população fechada, eventos corretamente pareados entre as fontes, e independência entre as fontes com probabilidade de captura homogênea. Com dependência positiva (um evento captado por uma fonte tem mais chance de ser captado pela outra) ou heterogeneidade, o total é subestimado; com dependência negativa, superestimado \parencite{hook1995}. Com duas fontes, a dependência não é identificável pelos dados, e o que se pode fazer é estratificar por fatores que afetam a captura, usar estimadores para captura heterogênea \parencite{chao1987} ou apresentar o resultado como limite sob hipóteses explícitas. O pareamento de eventos entre as fontes é um problema de ligação de registros \parencite{fellegi1969}, e erros de ligação viram erros de $m$.
}
\en{%
and $m/n_2$ estimates the capture probability of source 1. The estimator assumes a closed population, events correctly matched across sources, and independence between sources with homogeneous capture probability. With positive dependence (an event captured by one source is more likely to be captured by the other) or heterogeneity, the total is underestimated; with negative dependence, overestimated \parencite{hook1995}. With two sources, dependence is not identifiable from the data, and what one can do is stratify by factors that affect capture, use estimators for heterogeneous capture \parencite{chao1987} or present the result as a bound under explicit assumptions. Matching events across sources is a record linkage problem \parencite{fellegi1969}, and linkage errors become errors in $m$.
}

\section{\tr{Avaliação de modelos e previsão probabilística}{Model evaluation and probabilistic forecasting}}
\label{sec:fund-avaliacao}

\pt{%
Um modelo de risco produz, para cada veículo e mês, uma distribuição de probabilidade do número de acidentes, não um número. A avaliação adequada é por regras de pontuação próprias, que só são otimizadas em média pela distribuição verdadeira \parencite{gneiting2007}. A mais usada para contagens é o escore logarítmico, $\log p(y)$, a log-probabilidade que o modelo atribuiu ao valor observado. Para contagens, a calibração é verificada pela transformada integral de probabilidade não aleatorizada e por diagramas de calibração marginal \parencite{czado2009}. No ajuste, os resíduos quantílicos aleatorizados \parencite{dunn1996}, na versão simulada para modelos mistos do pacote DHARMa \parencite{hartig2026}, têm distribuição uniforme sob o modelo correto e permitem testar superdispersão e excesso de zeros.
}
\en{%
A risk model produces, for each vehicle and month, a probability distribution for the number of accidents, not a number. Proper evaluation uses proper scoring rules, which are optimized in expectation only by the true distribution \parencite{gneiting2007}. The most widely used one for counts is the log score, $\log p(y)$, the log probability the model assigned to the observed value. For counts, calibration is checked with the nonrandomized probability integral transform and marginal calibration diagrams \parencite{czado2009}. In fitting, randomized quantile residuals \parencite{dunn1996}, in the simulated version for mixed models of the DHARMa package \parencite{hartig2026}, are uniformly distributed under the correct model and allow testing for overdispersion and excess zeros.
}

\pt{%
A comparação de dois modelos pelo escore médio fora da amostra precisa considerar que as diferenças de escore são correlacionadas no tempo e dentro do mesmo veículo; o teste de \textcite{diebold1995} formaliza a comparação de acurácia preditiva para séries, e aqui é adaptado com bootstrap por veículo. A avaliação fora da amostra em séries temporais usa origem móvel: o modelo é reajustado a cada ponto de corte e avaliado nos períodos seguintes \parencite{tashman2000}.
}
\en{%
Comparing two models by their mean out-of-sample score must account for score differences being correlated over time and within the same vehicle; the \textcite{diebold1995} test formalizes the comparison of predictive accuracy for series, and here it is adapted with a bootstrap by vehicle. Out-of-sample evaluation of time series uses a rolling origin: the model is refitted at each cutoff and evaluated on the following periods \parencite{tashman2000}.
}

\pt{%
Para a extensão de previsão, as horas futuras são previstas como distribuição com métodos de séries temporais \parencite{hyndman2021}. O alisamento exponencial com tendência amortecida, que evita extrapolar tendências que tendem a não persistir, é um ponto de partida robusto \parencite{gardner1985}. A avaliação usa o CRPS e a perda quantílica \parencite{gneiting2007}.
}
\en{%
For the forecasting extension, future hours are forecast as a distribution with time series methods \parencite{hyndman2021}. Exponential smoothing with a damped trend, which avoids extrapolating trends that tend not to persist, is a robust starting point \parencite{gardner1985}. Evaluation uses the CRPS and the quantile loss \parencite{gneiting2007}.
}

\section{\tr{Testes de equivalência e testes múltiplos}{Equivalence and multiple tests}}
\label{sec:fund-testes}

\pt{%
Um teste de significância que não rejeita $H_0$ não mostra que o efeito é nulo: pode ter faltado poder. Para afirmar que um efeito é desprezível, inverte-se a hipótese nula. No procedimento de dois testes unilaterais \parencite{schuirmann1987}, fixa-se uma margem $(\delta_L,\ \delta_U)$ de efeitos considerados irrelevantes e testa-se $H_0\colon \theta \le \delta_L$ ou $\theta \ge \delta_U$ contra $H_1\colon \delta_L < \theta < \delta_U$. Rejeitar os dois testes unilaterais ao nível $\alpha$ equivale a o intervalo de confiança de $1 - 2\alpha$ caber inteiro dentro da margem; com $\alpha = 0{,}05$, usa-se o intervalo de 90\% \parencite{lakens2017}. A margem precisa ser fixada antes de ver os resultados finais, com justificativa substantiva.
}
\en{%
A significance test that fails to reject $H_0$ does not show that the effect is null: power may have been lacking. To claim that an effect is negligible, the null hypothesis is inverted. In the two one-sided tests procedure \parencite{schuirmann1987}, a margin $(\delta_L,\ \delta_U)$ of effects deemed irrelevant is fixed and one tests $H_0\colon \theta \le \delta_L$ or $\theta \ge \delta_U$ against $H_1\colon \delta_L < \theta < \delta_U$. Rejecting both one-sided tests at level $\alpha$ is equivalent to the $1 - 2\alpha$ confidence interval lying entirely within the margin; with $\alpha = 0.05$, the 90\% interval is used \parencite{lakens2017}. The margin must be fixed before seeing the final results, with a substantive justification.
}

\pt{%
Quando vários tipos de evento de condução são testados, a chance de ao menos um falso positivo cresce com o número de testes. O método de \textcite{holm1979} ordena os valores-$p$ e compara o $k$-ésimo menor a $\alpha/(m - k + 1)$, controlando a taxa de erro da família com mais poder que a correção de Bonferroni.
}
\en{%
When several driving event types are tested, the chance of at least one false positive grows with the number of tests. The \textcite{holm1979} method orders the $p$-values and compares the $k$-th smallest to $\alpha/(m - k + 1)$, controlling the familywise error rate with more power than the Bonferroni correction.
}

\section{\tr{Síntese e lacuna}{Summary and gap}}
\label{sec:fund-lacuna}

\pt{%
Os métodos necessários existem e são maduros: modelos de contagem com exposição e efeitos aleatórios, Bayes empírico e gráficos de funil, captura e recaptura, regras de pontuação e testes de equivalência. O que falta é aplicá-los juntos ao pátio. A literatura de logística trata do pátio pela eficiência; a de segurança de caminhões, da estrada; a de segurança em armazéns, das empilhadeiras. O caminhão de manobra, que anda pouco, manobra muito e tem acidentes de baixa velocidade sujeitos a subnotificação, não tem estudo próprio. Este trabalho testa se o que funciona na estrada (telemetria como indicador de risco, detecção automática de colisão) funciona no pátio, e propõe uma forma defensável de comparar o risco entre unidades.
}
\en{%
The required methods exist and are mature: count models with exposure and random effects, empirical Bayes and funnel plots, capture-recapture, scoring rules and equivalence tests. What is missing is applying them together to the yard. The logistics literature treats the yard from the efficiency side; truck safety research, the road; warehouse safety research, forklifts. The yard truck, which travels little, maneuvers a lot and has low-speed accidents subject to underreporting, has no study of its own. This work tests whether what works on the road (telematics as a risk indicator, automatic crash detection) works in the yard, and proposes a defensible way to compare risk across sites.
}

\chapter{\tr{Dados}{Data}}
\label{cap:dados}

\section{\tr{Dados da operação}{Operational data}}
\label{sec:dados-operacao}

\pt{%
O estudo usa apenas dados que a empresa já coleta na rotina; não há coleta própria. A Tabela~\ref{tab:fontes} resume as quatro fontes, com ordens de grandeza de um levantamento preliminar \preencher{atualizar período, cobertura e escala depois da autorização formal da empresa; até lá, só ordens de grandeza}.
}
\en{%
The study uses only data the company already collects routinely; there is no new data collection. Table~\ref{tab:fontes} summarizes the four sources, with orders of magnitude from a preliminary survey \preencher{update period, coverage and scale after the company's formal authorization; until then, orders of magnitude only}.
}

\begin{table}[ht]
\centering
\small
\caption{\tr{Fontes de dados da operação}{Operational data sources}}
\label{tab:fontes}
\begin{tabularx}{\textwidth}{>{\raggedright\arraybackslash}p{3.0cm} >{\raggedright\arraybackslash}X >{\raggedright\arraybackslash}p{3.6cm}}
\toprule
\tr{Fonte}{Source} & \tr{Conteúdo}{Content} & \tr{Escala aproximada}{Approximate scale} \\
\midrule
\tr{Relatos de acidentes}{Accident reports} &
\tr{Data, tipo de evento, unidade, número do veículo, gravidade, descrição, situação da revisão}%
   {Date, event type, site, vehicle number, severity, description, review status} &
\tr{Centenas de relatos por ano}{Hundreds of reports per year} \\
\tr{Telemetria veicular}{Vehicle telematics} &
\tr{Viagens (distância, duração), eventos de condução por tipo, colisões detectadas e resultado da revisão, unidade do veículo}%
   {Trips (distance, duration), driving events by type, detected collisions and review outcome, vehicle site} &
\tr{Algumas centenas de veículos; mais de um milhão de viagens}{A few hundred vehicles; over one million trips} \\
\tr{Horas trabalhadas}{Hours worked} &
\tr{Horas por unidade, departamento e mês, com horas extras}%
   {Hours per site, department and month, with overtime} &
\tr{Dezenas de unidades}{Dozens of sites} \\
\tr{Manutenção}{Maintenance} &
\tr{Ordens de reparo por veículo e data, com motivo (dano, acidente, preventiva)}%
   {Repair orders by vehicle and date, with reason (damage, accident, preventive)} &
\tr{Centenas de ordens por dano}{Hundreds of damage orders} \\
\bottomrule
\end{tabularx}
\end{table}

\subsection{\tr{Relatos de acidentes}{Accident reports}}

\pt{%
Os relatos são registrados por supervisores e motoristas num aplicativo próprio e revisados pela área de segurança. Cada relato traz data e hora, unidade, número do caminhão, tipo de evento (colisão com carreta, doca, estrutura fixa, outro veículo; lesão; quase-acidente), gravidade e uma descrição em texto livre. Antes do aplicativo, os relatos eram feitos por formulário eletrônico, com campos menos padronizados. O aplicativo teve um período de implantação em que relatos de teste se misturaram aos reais, e parte dos relatos fica pendente de revisão por algum tempo.
}
\en{%
Reports are recorded by supervisors and drivers in an in-house application and reviewed by the safety team. Each report carries date and time, site, truck number, event type (collision with trailer, dock, fixed structure, other vehicle; injury; near miss), severity and a free-text description. Before the application, reports were filed through an electronic form with less standardized fields. The application had a roll-out period in which test reports were mixed with real ones, and some reports stay pending review for a while.
}

\subsection{\tr{Telemetria}{Telematics}}

\pt{%
Os caminhões têm um dispositivo de telemetria comercial que registra cada viagem (início, fim, distância, duração), eventos de condução detectados por limiares de aceleração, velocidade e regras de trânsito, e colisões detectadas automaticamente por pico de aceleração. Cada detecção de colisão passa por uma revisão que a confirma ou descarta como falso alarme. O cadastro associa cada veículo a uma unidade por meio de etiquetas que usam o mesmo código de unidade da folha de pagamento. A sincronização desses dados com o ambiente analítico da empresa foi interrompida e precisa ser retomada para cobrir o período completo do projeto \preencher{confirmar o período efetivamente disponível}.
}
\en{%
The trucks carry a commercial telematics device that records each trip (start, end, distance, duration), driving events detected by acceleration, speed and traffic-rule thresholds, and collisions detected automatically by acceleration peaks. Each crash detection goes through a review that confirms it or dismisses it as a false alarm. The registry assigns each vehicle to a site through tags that use the same site code as payroll. Syncing these data into the company's analytics environment was paused and must be resumed to cover the project's full period \preencher{confirm the period actually available}.
}

\subsection{\tr{Horas trabalhadas e manutenção}{Hours worked and maintenance}}

\pt{%
A folha de pagamento fornece horas trabalhadas por unidade e mês, inclusive horas extras, sem identificar pessoas na extração usada. Ela é a medida de exposição do plano alternativo (Seção~\ref{sec:met-alternativo}) e uma segunda medida de exposição para comparar com o tempo de viagem da telemetria. O sistema de manutenção registra ordens de reparo por veículo e data, com o motivo; as ordens abertas por dano ou acidente são a segunda fonte para a subnotificação (Seção~\ref{sec:met-subnotificacao}) \preencher{verificar se o motivo ``acidente'' é atribuído só depois do relato, o que tornaria as fontes dependentes por construção}.
}
\en{%
Payroll provides hours worked per site and month, including overtime, without identifying people in the extraction used. It is the exposure measure of the fallback plan (Section~\ref{sec:met-alternativo}) and a second exposure measure to compare with telematics trip time. The maintenance system records repair orders by vehicle and date, with the reason; orders opened for damage or accident are the second source for underreporting (Section~\ref{sec:met-subnotificacao}) \preencher{check whether the ``accident'' reason is assigned only after the report, which would make the sources dependent by construction}.
}

\subsection{\tr{Definições}{Definitions}}
\label{sec:definicoes}

\pt{%
\begin{description}[leftmargin=0.6cm,style=nextline]
  \item[\textnormal{\textit{Unidade de análise.}}] Veículo $\times$ mês. A unidade (site) de cada veículo é a do cadastro da telemetria no mês; veículos que mudam de unidade contribuem para mais de uma.
  \item[\textnormal{\textit{Desfecho principal.}}] Número de colisões relatadas no mês, com ou sem lesão, depois de excluir relatos de teste e duplicados. Lesões e quase-acidentes entram como definições alternativas em análise de sensibilidade.
  \item[\textnormal{\textit{Exposição principal.}}] Horas de viagem registradas pela telemetria. A distância percorrida é a exposição alternativa.
  \item[\textnormal{\textit{Tipo de uso.}}] Manobra no pátio (\textit{spotting}), transferência entre instalações (\textit{shuttling}) e uso local, atribuído a cada veículo pelo padrão de viagens e conferido com o cadastro \preencher{documentar a regra de atribuição}.
  \item[\textnormal{\textit{Eventos de condução.}}] Taxas de cada tipo de evento por hora de viagem no mês, padronizadas (média zero e desvio padrão um) para que as razões de taxas sejam comparáveis entre tipos \preencher{listar os tipos de evento usados}.
\end{description}
}
\en{%
\begin{description}[leftmargin=0.6cm,style=nextline]
  \item[\textnormal{\textit{Unit of analysis.}}] Vehicle $\times$ month. Each vehicle's site is the one in the telematics registry for that month; vehicles that move between sites contribute to more than one.
  \item[\textnormal{\textit{Main outcome.}}] Number of reported collisions in the month, with or without injury, after excluding test and duplicate reports. Injuries and near misses enter as alternative definitions in a sensitivity analysis.
  \item[\textnormal{\textit{Main exposure.}}] Trip hours recorded by telematics. Distance traveled is the alternative exposure.
  \item[\textnormal{\textit{Use type.}}] Yard spotting, shuttling between facilities and local use, assigned to each vehicle by its trip pattern and checked against the registry \preencher{document the assignment rule}.
  \item[\textnormal{\textit{Driving events.}}] Rates of each event type per trip hour in the month, standardized (zero mean and unit standard deviation) so that rate ratios are comparable across types \preencher{list the event types used}.
\end{description}
}

\subsection{\tr{Ligação entre as fontes e qualidade}{Linkage across sources and quality}}

\pt{%
O número do caminhão informado nos relatos coincide com o identificador da telemetria na maior parte dos casos depois de uma normalização simples (caixa, espaços, prefixos), e o restante é tratado por comparação aproximada (Seção~\ref{sec:met-painel}). Os problemas de qualidade já conhecidos são: relatos de teste no período de implantação; unidade e veículo digitados em texto livre em parte dos relatos; viagens com registro de fim inválido; veículos sem unidade fixa, como os da oficina central; e mudanças de veículos entre unidades. Cada regra de limpeza é documentada, e seu efeito sobre os resultados é testado em análise de sensibilidade.
}
\en{%
The truck number given in reports matches the telematics identifier in most cases after simple normalization (case, spaces, prefixes), and the rest is handled by approximate matching (Section~\ref{sec:met-painel}). Known quality issues are: test reports during the roll-out period; site and vehicle typed as free text in some reports; trips with invalid end records; vehicles without a fixed site, such as those of the central workshop; and vehicles moving between sites. Every cleaning rule is documented, and its effect on results is tested in a sensitivity analysis.
}

\begin{nota}[\tr{Privacidade e confidencialidade}{Privacy and confidentiality}]
\tr%
  {Nenhum dado identifica pessoas. A análise usa veículo, unidade e mês; identificadores de motorista, nomes, imagens de câmera e trajetos de GPS não são usados. Os dados não saem do ambiente da empresa; só tabelas agregadas entram no repositório da pesquisa e no texto, que não nomeia a empresa, seus clientes nem suas unidades. As unidades aparecem mascaradas (Unidade 01, 02...) e os resultados são apresentados como valores relativos e ordens de grandeza.}%
  {No data identifies people. The analysis uses vehicle, site and month; driver identifiers, names, camera footage and GPS tracks are not used. Data stays within the company's environment; only aggregated tables enter the research repository and the text, which does not name the company, its clients or its sites. Sites appear masked (Site 01, 02...) and results are presented as relative values and orders of magnitude.}
\end{nota}

\section{\tr{Dados públicos usados nos estudos preliminares}{Public data used in the preliminary studies}}
\label{sec:dados-publicos}

\pt{%
O conjunto \textit{Logistics Operations Database} \parencite{yogape2025}, publicado no Kaggle sob licença MIT, simula uma transportadora fictícia dos EUA entre 2022 e 2024, com 14 tabelas ligadas por chaves: caminhões, terminais, viagens, cargas e incidentes de segurança, além de uma tabela mensal por caminhão com milhas rodadas. Ele é usado aqui só como \emph{estrutura realista} (quantos caminhões, quanto cada um roda por mês, como o volume varia por semana), porque os incidentes do conjunto se mostraram praticamente independentes de caminhão, terminal e exposição (\texttt{notas/05-datasets-publicos.md}). Os acidentes usados nos estudos de poder são, portanto, simulados sobre essa estrutura, com efeitos conhecidos (Seção~\ref{sec:poder}).
}
\en{%
The \textit{Logistics Operations Database} \parencite{yogape2025}, published on Kaggle under the MIT license, simulates a fictitious US trucking company from 2022 to 2024, with 14 linked tables: trucks, terminals, trips, loads and safety incidents, plus a monthly truck table with miles driven. It is used here only as a \emph{realistic structure} (how many trucks, how much each drives per month, how volume varies by week), because the incidents in the dataset turned out to be practically independent of truck, terminal and exposure. The accidents used in the power studies are therefore simulated on this structure, with known effects (Section~\ref{sec:poder}).
}

\chapter{\tr{Análise descritiva}{Descriptive analysis}}
\label{cap:descritiva}

\pt{%
Este capítulo descreve as métricas que os modelos dos capítulos seguintes vão usar: a exposição, o desfecho, a concentração dos acidentes entre veículos e unidades, o padrão no tempo, os tipos e causas de acidente, os eventos de telemetria e as fontes disponíveis para estimar a subnotificação. O objetivo é duplo: conhecer a distribuição de cada métrica antes de modelá-la e justificar, com os dados, as escolhas do Capítulo~\ref{cap:metodos}. A análise preditiva (Capítulo~\ref{cap:preditiva}) e a prescritiva (Capítulo~\ref{cap:prescritiva}) partem daqui.
}
\en{%
This chapter describes the metrics that the models in the following chapters will use: exposure, the outcome, the concentration of accidents across vehicles and sites, the pattern over time, accident types and causes, telematics events and the sources available to estimate underreporting. The aim is twofold: to know the distribution of each metric before modeling it and to justify, with the data, the choices in Chapter~\ref{cap:metodos}. The predictive (Chapter~\ref{cap:preditiva}) and prescriptive (Chapter~\ref{cap:prescritiva}) analyses build on it.
}

\section{\tr{Convenções}{Conventions}}

\pt{%
Os dados são os da operação descritos no Capítulo~\ref{cap:dados}, extraídos só como agregados por veículo e mês ou por categoria, sem dados pessoais. O painel cobre doze meses comuns à telemetria e aos relatos e inclui os veículo-meses com ao menos dez horas de viagem; os relatos e as fontes auxiliares cobrem o período disponível de cada uma. Para preservar a confidencialidade, todas as figuras e números do capítulo são relativos: proporções, índices em relação à média, razões e exposição em relação à mediana. As unidades não são identificadas. As figuras são geradas por \texttt{analise/real/descritiva.py} (Python) e \texttt{r/descritiva.R} (R), e a Tabela~\ref{tab:descritiva} resume os números citados.
}
\en{%
The data are the operational data described in Chapter~\ref{cap:dados}, extracted only as aggregates per vehicle and month or per category, with no personal data. The panel covers twelve months common to telematics and reports and includes vehicle-months with at least ten trip hours; reports and auxiliary sources cover each one's available period. To preserve confidentiality, all figures and numbers in the chapter are relative: proportions, indices relative to the mean, ratios and exposure relative to the median. Sites are not identified. Figures are generated by \texttt{analise/real/descritiva.py} (Python) and \texttt{r/descritiva.R} (R), and Table~\ref{tab:descritiva} summarizes the numbers cited.
}

\section{\tr{Exposição}{Exposure}}
\label{sec:desc-exposicao}

\begin{figure}[!htbp]
\centering
\includegraphics[width=\textwidth]{\arqfig{fig_desc_exposicao}}
\caption{\tr%
  {Exposição. (a) a (c): distribuição por veículo-mês, por tipo de uso, relativa à mediana geral (escala log; cada ponto é um veículo-mês). (d) Fração das horas pagas que são horas extras, por unidade e mês da folha; linha tracejada: mediana.}%
  {Exposure. (a) to (c): distribution per vehicle-month, by use type, relative to the overall median (log scale; each point is a vehicle-month). (d) Share of paid hours that are overtime, per payroll site and month; dashed line: median.}}
\label{fig:desc-exposicao}
\end{figure}

\pt{%
As horas de viagem por veículo-mês têm distribuição parecida nos três tipos de uso, com mediana próxima da geral e dispersão grande dentro de cada tipo (Figura~\ref{fig:desc-exposicao}a). As milhas, ao contrário, separam os tipos por uma ordem de grandeza (Figura~\ref{fig:desc-exposicao}b): o spotter anda a cerca de \descVelSpotter{} da velocidade média do veículo local, e o shuttle, a cerca de \descVelShuttle{} (Figura~\ref{fig:desc-exposicao}c). É a forma visual do argumento do Capítulo~\ref{cap:metodos}: usar a distância como exposição faria o spotter parecer muito mais arriscado só porque quase não roda, e por isso a exposição principal é o tempo. Na folha de pagamento, as horas extras representam em mediana cerca de \descHorasExtras\% das horas de cada unidade-mês, com muitas unidades-mês sem nenhuma e uma cauda longa (Figura~\ref{fig:desc-exposicao}d). Essa variação é o que torna a fração de horas extras uma covariável candidata \parencite{dembe2005}.
}
\en{%
Trip hours per vehicle-month have a similar distribution across the three use types, with a median close to the overall one and a large spread within each type (Figure~\ref{fig:desc-exposicao}a). Miles, in contrast, separate the types by an order of magnitude (Figure~\ref{fig:desc-exposicao}b): the spotter moves at about \descVelSpotter{} of the local vehicle's average speed, and the shuttle at about \descVelShuttle{} (Figure~\ref{fig:desc-exposicao}c). This is the visual form of the argument in Chapter~\ref{cap:metodos}: using distance as exposure would make the spotter look much riskier only because it barely moves, which is why the main exposure is time. In payroll, overtime accounts for a median of about \descHorasExtras\% of each site-month's hours, with many site-months with none and a long tail (Figure~\ref{fig:desc-exposicao}d). This variation is what makes the overtime share a candidate covariate \parencite{dembe2005}.
}

\section{\tr{Desfecho: contagens e superdispersão}{Outcome: counts and overdispersion}}
\label{sec:desc-desfecho}

\begin{figure}[!htbp]
\centering
\includegraphics[width=\textwidth]{\arqfig{fig_desc_contagem}}
\caption{\tr%
  {Desfecho. (a) Proporção de veículo-meses com 0, 1, 2 e 3 ou mais acidentes (barras) e o esperado por três modelos com tipo de uso e horas de viagem como \textit{offset} (pontos; média de 300 simulações). (b) O mesmo para o total de acidentes de cada veículo no período. (c) Risco relativo estimado de cada veículo (efeito aleatório da binomial negativa), com intervalo de 95\%. Escala raiz em (a) e (b), log em (c).}%
  {Outcome. (a) Share of vehicle-months with 0, 1, 2 and 3 or more accidents (bars) and the expected share under three models with use type and trip hours as offset (points; mean of 300 simulations). (b) The same for each vehicle's total accidents in the period. (c) Estimated relative risk of each vehicle (negative binomial random effect), with 95\% interval. Square-root scale in (a) and (b), log in (c).}}
\label{fig:desc-contagem}
\end{figure}

\pt{%
O acidente é raro no nível do veículo-mês: cerca de \descZeros\% dos veículo-meses não têm nenhum, e a variância das contagens é apenas \descVarMedia{} vez a média. Nesse nível, o próprio Poisson com exposição e tipo de uso já reproduz a proporção de zeros (\descRzerosPoisson\% esperados contra \descRzerosObs\% observados; Figura~\ref{fig:desc-contagem}a), o que confirma que não há excesso de zeros a modelar \parencite{lord2005}. A heterogeneidade aparece quando se soma o período de cada veículo (Figura~\ref{fig:desc-contagem}b): o Poisson espera \descRveicZeroPoisson\% de veículos sem acidente, contra \descRveicZeroObs\% observados, e a binomial negativa com efeito de veículo acerta essa proporção (\descRveicZeroRE\%). O desvio padrão do efeito de veículo é de \descRdpVeiculo{} na escala log, mas os intervalos de cada veículo são largos e quase todos cruzam 1 (Figura~\ref{fig:desc-contagem}c): com um ano de dados, dá para afirmar que os veículos diferem, mas não para dizer com segurança quais são os piores. É o cenário em que o encolhimento é necessário.
}
\en{%
Accidents are rare at the vehicle-month level: about \descZeros\% of vehicle-months have none, and the variance of the counts is only \descVarMedia{} times the mean. At this level, the Poisson model with exposure and use type already reproduces the share of zeros (\descRzerosPoisson\% expected against \descRzerosObs\% observed; Figure~\ref{fig:desc-contagem}a), which confirms there is no excess of zeros to model \parencite{lord2005}. Heterogeneity appears when each vehicle's period is summed (Figure~\ref{fig:desc-contagem}b): the Poisson model expects \descRveicZeroPoisson\% of vehicles with no accident, against \descRveicZeroObs\% observed, and the negative binomial with a vehicle effect matches this share (\descRveicZeroRE\%). The standard deviation of the vehicle effect is \descRdpVeiculo{} on the log scale, but each vehicle's interval is wide and almost all cross 1 (Figure~\ref{fig:desc-contagem}c): with one year of data, one can claim that vehicles differ, but not say with confidence which ones are the worst. This is the setting in which shrinkage is needed.
}

\section{\tr{Concentração e diferenças entre unidades}{Concentration and differences across sites}}
\label{sec:desc-concentracao}

\begin{figure}[!htbp]
\centering
\includegraphics[width=\textwidth]{\arqfig{fig_desc_concentracao}}
\caption{\tr%
  {Concentração. (a) e (b) Curvas de concentração: fração acumulada dos acidentes (linha cheia) e das horas (tracejada) contra a fração acumulada de veículos ou unidades, ordenados do maior para o menor; faixa: 95\% das curvas simuladas sob Poisson com taxa constante dentro de cada tipo de uso. (c) Razão entre acidentes observados e esperados por unidade (mesmo modelo), com intervalo exato de 95\%, para as unidades com ao menos um acidente esperado; \corpb{em vermelho}{em preto}, as que excluem 1. (d) Horas de viagem por unidade relativas à mediana.}%
  {Concentration. (a) and (b) Concentration curves: cumulative share of accidents (solid line) and hours (dashed) against the cumulative share of vehicles or sites, ordered from largest to smallest; band: 95\% of curves simulated under Poisson with a constant rate within each use type. (c) Ratio of observed to expected accidents per site (same model), with exact 95\% interval, for sites with at least one expected accident; \corpb{in red}{in black}, those that exclude 1. (d) Trip hours per site relative to the median.}}
\label{fig:desc-concentracao}
\end{figure}

\pt{%
Os acidentes parecem muito concentrados: metade dos veículos não teve nenhum, e o coeficiente de Gini dos acidentes entre veículos é de \descGiniObs{} (Figura~\ref{fig:desc-concentracao}a). Mas quase toda essa concentração é a esperada por acaso e exposição: um Poisson com taxa constante dentro de cada tipo de uso produz Gini de \descGiniPoisson{}, e a curva observada fica só um pouco acima da faixa simulada. Entre unidades, o padrão se repete (Gini de \descGiniSiteObs{} contra \descGiniSitePoisson{}; Figura~\ref{fig:desc-concentracao}b). A consequência prática é importante para a gestão: ver que poucos veículos ou unidades concentram os acidentes não é, por si só, evidência de que eles sejam mais arriscados.
}
\en{%
Accidents look highly concentrated: half of the vehicles had none, and the Gini coefficient of accidents across vehicles is \descGiniObs{} (Figure~\ref{fig:desc-concentracao}a). But almost all of this concentration is what chance and exposure would produce: a Poisson model with a constant rate within each use type yields a Gini of \descGiniPoisson{}, and the observed curve lies only slightly above the simulated band. Across sites the pattern repeats (Gini of \descGiniSiteObs{} against \descGiniSitePoisson{}; Figure~\ref{fig:desc-concentracao}b). The practical consequence matters for management: seeing that a few vehicles or sites concentrate accidents is not, in itself, evidence that they are riskier.
}

\pt{%
As unidades diferem muito em tamanho: a que está no percentil 90 da exposição opera cerca de \descPnoventaPdez{} vezes mais horas que a do percentil 10 (Figura~\ref{fig:desc-concentracao}d). Por isso a contagem bruta mede sobretudo o tamanho. Ajustando pela exposição e pelo tipo de uso, cerca de \descForaIC\% das unidades com ao menos um acidente esperado ficam com intervalo de 95\% inteiramente acima ou abaixo de 1 (Figura~\ref{fig:desc-concentracao}c). A proporção é menor que a da exploração do Capítulo~\ref{cap:viabilidade}, que não ajustava pelo tipo de uso: parte da diferença aparente entre unidades vem da mistura de operações de cada uma. O que sobra é a variação que o efeito de unidade do modelo hierárquico deve capturar (H1b).
}
\en{%
Sites differ greatly in size: the one at the 90th percentile of exposure operates about \descPnoventaPdez{} times more hours than the one at the 10th percentile (Figure~\ref{fig:desc-concentracao}d). That is why the raw count mostly measures size. Adjusting for exposure and use type, about \descForaIC\% of sites with at least one expected accident have a 95\% interval entirely above or below 1 (Figure~\ref{fig:desc-concentracao}c). The share is smaller than in the exploration of Chapter~\ref{cap:viabilidade}, which did not adjust for use type: part of the apparent difference across sites comes from each site's mix of operations. What remains is the variation the hierarchical model's site effect should capture (H1b).
}

\section{\tr{Padrão no tempo}{Pattern over time}}
\label{sec:desc-tempo}

\begin{figure}[!htbp]
\centering
\includegraphics[width=\textwidth]{\arqfig{fig_desc_tempo}}
\caption{\tr%
  {Tempo. (a) Relatos de acidente por mês relativos à média mensal; hachurado: relatos pendentes de revisão no período de implantação do aplicativo, suspeitos de teste. (b) e (c) Distribuição dos acidentes (sem os suspeitos) por hora do dia e por dia da semana; linha pontilhada: distribuição uniforme; faixas: turnos de 8 horas. (d) Horas da folha por semana e horas de viagem por mês, relativas à média de cada série.}%
  {Time. (a) Accident reports per month relative to the monthly mean; hatched: reports pending review during the application roll-out, suspected tests. (b) and (c) Distribution of accidents (excluding suspected tests) by hour of day and day of week; dotted line: uniform distribution; bands: 8-hour shifts. (d) Payroll hours per week and trip hours per month, relative to each series' mean.}}
\label{fig:desc-tempo}
\end{figure}

\pt{%
O número de relatos varia bastante de um mês para outro (Figura~\ref{fig:desc-tempo}a), e cerca de \descSuspeitos\% dos relatos de acidente são suspeitos de teste: estavam pendentes de revisão e foram enviados durante a implantação do aplicativo, concentrados entre maio e outubro de 2026. A regra de exclusão desses relatos muda o desfecho de forma desigual no tempo e precisa de análise de sensibilidade, como previsto no Capítulo~\ref{cap:metodos}.
}
\en{%
The number of reports varies considerably from one month to the next (Figure~\ref{fig:desc-tempo}a), and about \descSuspeitos\% of accident reports are suspected tests: they were pending review and submitted during the application roll-out, concentrated between May and October 2026. The rule for excluding these reports changes the outcome unevenly over time and needs a sensitivity analysis, as planned in Chapter~\ref{cap:metodos}.
}

\pt{%
Quase todos os relatos (\descComHora\%) trazem a hora do acidente. Os acidentes se distribuem pelas 24 horas, com \descTurnoUm\%, \descTurnoDois\% e \descTurnoTres\% nos turnos da manhã, da tarde e da noite, e um pico no início da noite (Figura~\ref{fig:desc-tempo}b). Por dia da semana, há mais acidentes de terça a quinta e menos no fim de semana (Figura~\ref{fig:desc-tempo}c), o que acompanha o volume de operação. Sem a exposição por hora e por dia, essas distribuições ainda não podem ser lidas como risco; servem para mostrar que um modelo por turno é viável. Já a exposição cresceu de forma contínua ao longo do ano, tanto nas horas da folha quanto nas horas de viagem (Figura~\ref{fig:desc-tempo}d), com a entrada de novas unidades e veículos: comparar meses sem \textit{offset} e sem efeito de calendário confundiria crescimento da operação com aumento de risco.
}
\en{%
Almost all reports (\descComHora\%) include the time of the accident. Accidents are spread over the 24 hours, with \descTurnoUm\%, \descTurnoDois\% and \descTurnoTres\% in the morning, afternoon and night shifts, and a peak in the early evening (Figure~\ref{fig:desc-tempo}b). By day of week, there are more accidents from Tuesday to Thursday and fewer on weekends (Figure~\ref{fig:desc-tempo}c), which follows operating volume. Without exposure by hour and by day, these distributions cannot yet be read as risk; they show that a shift-level model is feasible. Exposure, in turn, grew steadily over the year, in both payroll hours and trip hours (Figure~\ref{fig:desc-tempo}d), as new sites and vehicles came in: comparing months without an offset and a calendar effect would confuse operational growth with rising risk.
}

\section{\tr{Tipos e causas de acidente}{Accident types and causes}}
\label{sec:desc-tipos}

\begin{figure}[!htbp]
\centering
\includegraphics[width=\textwidth]{\arqfig{fig_desc_tipos}}
\caption{\tr%
  {Tipos e causas. (a) Distribuição dos relatos por tipo. (b) Participação de cada tipo de uso nos acidentes, nas horas de viagem e nas milhas. (c) Causa atribuída aos acidentes, em categorias fechadas do formulário de relato.}%
  {Types and causes. (a) Distribution of reports by type. (b) Share of each use type in accidents, trip hours and miles. (c) Cause assigned to accidents, in closed categories of the report form.}}
\label{fig:desc-tipos}
\end{figure}

\pt{%
Os acidentes com veículo são cerca de \descAcidentesPct\% dos relatos; lesões e outros eventos são minoria (Figura~\ref{fig:desc-tipos}a), o que reforça a escolha de todas as colisões como desfecho principal. O spotter responde pela maior parte dos acidentes, por pouco mais da metade das horas e por uma fração pequena das milhas; o veículo local, pelo inverso (Figura~\ref{fig:desc-tipos}b). As causas mais frequentes são as do acidente de baixa velocidade típico do pátio: problemas na porta da carreta e na fixação de equipamento, erro operacional e ré com folga mal avaliada somam cerca de \descTopCausas\% dos acidentes, seguidas de colisão com objeto fixo, carreta caída e problemas de engate (Figura~\ref{fig:desc-tipos}c). Colisões com outros veículos, o acidente típico da estrada, são minoria. É um retrato que antecipa por que eventos de telemetria desenhados para a estrada podem carregar pouco sinal no pátio (H2).
}
\en{%
Vehicle accidents are about \descAcidentesPct\% of reports; injuries and other events are a minority (Figure~\ref{fig:desc-tipos}a), which supports choosing all collisions as the main outcome. The spotter accounts for most accidents, a little over half of the hours and a small share of the miles; the local vehicle, for the opposite (Figure~\ref{fig:desc-tipos}b). The most frequent causes are those of the low-speed accident typical of the yard: trailer door and equipment securing problems, operational error and backing with misjudged clearance add up to about \descTopCausas\% of accidents, followed by striking a fixed object, dropped trailers and coupling problems (Figure~\ref{fig:desc-tipos}c). Collisions with other vehicles, the typical road accident, are a minority. This picture anticipates why telematics events designed for the road may carry little signal in the yard (H2).
}

\section{\tr{Eventos de telemetria}{Telematics events}}
\label{sec:desc-eventos}

\begin{figure}[!htbp]
\centering
\includegraphics[width=\textwidth]{\arqfig{fig_desc_eventos}}
\caption{\tr%
  {Eventos de telemetria. (a) Proporção de veículo-meses com ao menos um evento de cada tipo. (b) Correlação de Spearman entre as taxas de eventos por hora de viagem, a taxa de acidentes por hora e as horas de viagem, por veículo-mês.}%
  {Telematics events. (a) Share of vehicle-months with at least one event of each type. (b) Spearman correlation between event rates per trip hour, the accident rate per hour and trip hours, per vehicle-month.}}
\label{fig:desc-eventos}
\end{figure}

\pt{%
A maioria dos tipos de evento é rara: em boa parte dos veículo-meses não há nenhum, e só a ausência de cinto aparece na maioria deles (Figura~\ref{fig:desc-eventos}a). As taxas dos diferentes tipos são pouco correlacionadas entre si; a maior correlação, de \descRcorMax{}, é entre frenagem brusca e curva brusca, dois eventos cinemáticos. A correlação de cada tipo com a taxa de acidentes do mesmo veículo-mês não passa de \descRcorAcidMax{} em módulo, e com as horas de viagem, de \descRcorHorasMax{} (Figura~\ref{fig:desc-eventos}b). Duas consequências para o Capítulo~\ref{cap:preditiva}: testar cada tipo num modelo próprio não perde muito por ignorar os demais, e o sinal marginal da telemetria é fraco já antes de qualquer ajuste, o que reforça a necessidade dos testes de equivalência de H2b.
}
\en{%
Most event types are rare: many vehicle-months have none, and only the no-seatbelt event appears in most of them (Figure~\ref{fig:desc-eventos}a). The rates of different types are weakly correlated with each other; the largest correlation, \descRcorMax{}, is between harsh braking and harsh turns, two kinematic events. Each type's correlation with the accident rate of the same vehicle-month does not exceed \descRcorAcidMax{} in absolute value, and with trip hours, \descRcorHorasMax{} (Figure~\ref{fig:desc-eventos}b). Two consequences for Chapter~\ref{cap:preditiva}: testing each type in its own model loses little by ignoring the others, and the marginal telematics signal is weak even before any adjustment, which reinforces the need for the equivalence tests of H2b.
}

\section{\tr{Fontes para a subnotificação}{Sources for underreporting}}
\label{sec:desc-fontes}

\begin{figure}[!htbp]
\centering
\includegraphics[width=\textwidth]{\arqfig{fig_desc_fontes}}
\caption{\tr%
  {Fontes. (a) Relatos de acidente, ordens de manutenção por dano e colisões detectadas pela telemetria, por mês, relativos à média de cada série no período comum. (b) Ordens de manutenção por dano relativas à média mensal; tons: grafia antiga e nova do motivo do reparo; linha tracejada: mudança de regime. (c) Nota das auditorias de segurança por unidade (\% dos pontos possíveis); linha tracejada: mediana. (d) Notas por mês.}%
  {Sources. (a) Accident reports, damage repair orders and telematics-detected collisions per month, relative to each series' mean over the common period. (b) Damage repair orders relative to the monthly mean; shades: old and new coding of the repair reason; dashed line: regime change. (c) Site safety audit scores (\% of possible points); dashed line: median. (d) Scores by month.}}
\label{fig:desc-fontes}
\end{figure}

\pt{%
As três fontes de acidentes andam de forma diferente ao longo do ano (Figura~\ref{fig:desc-fontes}a). As colisões detectadas pela telemetria crescem quase continuamente, acompanhando a frota, enquanto os relatos e as ordens de dano oscilam sem seguir umas às outras mês a mês, o que já sugere que as fontes captam universos diferentes. As ordens de manutenção por dano mudam de regime em novembro de 2025 (Figura~\ref{fig:desc-fontes}b): a grafia do motivo muda, o volume passa a ser de duas a três vezes o anterior, e o motivo que indicava acidente já relatado praticamente desaparece. Para a Etapa~4, isso tem dois lados. No período comum com os relatos, a dúvida sobre esse motivo quase não pesa; em compensação, é preciso entender o que passou a ser registrado como dano antes de usar as ordens como segunda fonte \preencher{verificar com a área de manutenção o que mudou em novembro de 2025}.
}
\en{%
The three accident sources move differently over the year (Figure~\ref{fig:desc-fontes}a). Telematics-detected collisions grow almost continuously, following the fleet, while reports and damage orders fluctuate without tracking each other month by month, which already suggests that the sources capture different universes. Damage repair orders change regime in November 2025 (Figure~\ref{fig:desc-fontes}b): the coding of the reason changes, volume becomes two to three times the previous level, and the reason that flagged an already reported accident practically disappears. For Stage~4 this cuts both ways. In the period shared with reports, the doubt about that reason barely matters; on the other hand, one must understand what came to be recorded as damage before using the orders as a second source \preencher{check with the maintenance team what changed in November 2025}.
}

\pt{%
As auditorias de segurança por unidade, registradas desde meados de 2026, têm nota mediana de \descNotaAud\% dos pontos possíveis, com a maioria entre 90\% e 100\% (Figura~\ref{fig:desc-fontes}c e d). A nota é uma medida candidata do clima de segurança de cada unidade, que a literatura associa à propensão a relatar \parencite{probst2008,probst2010}, mas o efeito teto limita a variação disponível para explicar diferenças entre unidades.
}
\en{%
Site safety audits, recorded since mid-2026, have a median score of \descNotaAud\% of possible points, with most between 90\% and 100\% (Figure~\ref{fig:desc-fontes}c and d). The score is a candidate measure of each site's safety climate, which the literature associates with the propensity to report \parencite{probst2008,probst2010}, but the ceiling effect limits the variation available to explain differences across sites.
}

\section{\tr{Síntese}{Summary}}

\begin{table}[!htbp]
\centering\small
\caption{\tr{Resumo descritivo (só valores relativos).}{Descriptive summary (relative values only).}}
\label{tab:descritiva}
\begin{tabularx}{\textwidth}{X r}
\toprule
\tr{Métrica}{Metric} & \tr{Valor}{Value} \\
\midrule
\tabinput{tabelas/tab_descritiva.tex}
\bottomrule
\end{tabularx}
\end{table}

\pt{%
\begin{enumerate}
  \item \textbf{Exposição pelo tempo, ajuste pelo tipo de uso.} As milhas separam os tipos de uso por uma ordem de grandeza; as horas, não.
  \item \textbf{Binomial negativa com efeito de veículo, sem inflação de zeros.} O excesso de zeros não existe no veículo-mês; a heterogeneidade aparece entre veículos.
  \item \textbf{Concentração não é risco.} Quase toda a concentração dos acidentes em poucos veículos e unidades é a esperada por acaso e exposição; comparar exige taxas ajustadas e intervalos.
  \item \textbf{Calendário e limpeza.} A operação cresceu ao longo do ano, e um quarto dos relatos de acidente precisa de regra de exclusão documentada.
  \item \textbf{Turno é viável.} A hora do acidente está quase sempre preenchida, o que permite modelos por turno.
  \item \textbf{Telemetria com sinal marginal fraco.} Eventos raros, pouco correlacionados entre si e com os acidentes.
  \item \textbf{Fontes de subnotificação com ressalvas.} A manutenção mudou de regime, e as auditorias têm efeito teto.
\end{enumerate}
}
\en{%
\begin{enumerate}
  \item \textbf{Time-based exposure, adjustment for use type.} Miles separate use types by an order of magnitude; hours do not.
  \item \textbf{Negative binomial with a vehicle effect, no zero inflation.} There is no excess of zeros at the vehicle-month level; heterogeneity appears across vehicles.
  \item \textbf{Concentration is not risk.} Almost all the concentration of accidents in a few vehicles and sites is what chance and exposure produce; comparison requires adjusted rates and intervals.
  \item \textbf{Calendar and cleaning.} Operations grew over the year, and a quarter of accident reports needs a documented exclusion rule.
  \item \textbf{Shifts are feasible.} The accident time is almost always filled in, which allows shift-level models.
  \item \textbf{Weak marginal telematics signal.} Rare events, weakly correlated with each other and with accidents.
  \item \textbf{Underreporting sources with caveats.} Maintenance changed regime, and audits have a ceiling effect.
\end{enumerate}
}

\chapter{\tr{Métodos}{Methods}}
\label{cap:metodos}

\pt{%
O trabalho tem quatro etapas: a integração das fontes num painel (Seção~\ref{sec:met-painel}), os modelos de risco (Seções~\ref{sec:met-modelo} a \ref{sec:met-indicadores}), a validação fora da amostra e os testes das hipóteses (Seções~\ref{sec:met-validacao} e \ref{sec:testes}) e a estimação da subnotificação (Seção~\ref{sec:met-subnotificacao}). O capítulo termina com o plano alternativo, as análises de sensibilidade, a extensão de previsão e as ferramentas.
}
\en{%
The work has four stages: integrating the sources into a panel (Section~\ref{sec:met-painel}), the risk models (Sections~\ref{sec:met-modelo} to \ref{sec:met-indicadores}), out-of-sample validation and hypothesis tests (Sections~\ref{sec:met-validacao} and \ref{sec:testes}) and underreporting estimation (Section~\ref{sec:met-subnotificacao}). The chapter ends with the fallback plan, sensitivity analyses, the forecasting extension and tools.
}

\section{\tr{Etapa 1: integração e painel}{Stage 1: integration and panel}}
\label{sec:met-painel}

\pt{%
As fontes são integradas num painel com uma linha por veículo e mês, contendo: o número de colisões relatadas (desfecho), as horas de viagem e a distância (exposições), a contagem de cada tipo de evento de condução, o número de colisões detectadas pelos sensores e o resultado da revisão, o tipo de uso, a unidade do veículo no mês e, por unidade e mês, as horas trabalhadas e as horas extras da folha.
}
\en{%
The sources are integrated into a panel with one row per vehicle and month, containing: the number of reported collisions (outcome), trip hours and distance (exposures), the count of each driving event type, the number of sensor-detected collisions and the review outcome, the use type, the vehicle's site in the month and, per site and month, payroll hours worked and overtime.
}

\pt{%
A ligação dos relatos aos veículos tem duas fases. Na determinística, o número do caminhão informado no relato é normalizado (caixa, espaços, zeros à esquerda, prefixos da unidade) e comparado ao nome do veículo na telemetria. Na aproximada, os relatos que sobram são comparados aos veículos da mesma unidade por similaridade de texto do identificador, com a regra de decisão inspirada em \textcite{fellegi1969}: um par é aceito quando a similaridade passa de um limiar e não há outro candidato próximo. A qualidade é medida numa amostra aleatória de relatos revisada manualmente, com precisão (pares aceitos corretos) e cobertura (pares corretos encontrados). Os modelos principais são reajustados só com as ligações determinísticas, como teste de sensibilidade.
}
\en{%
Linking reports to vehicles has two phases. In the deterministic phase, the truck number given in the report is normalized (case, spaces, leading zeros, site prefixes) and compared with the vehicle name in telematics. In the approximate phase, the remaining reports are compared with vehicles at the same site by text similarity of the identifier, with a decision rule inspired by \textcite{fellegi1969}: a pair is accepted when similarity exceeds a threshold and there is no other close candidate. Quality is measured on a random sample of reports reviewed by hand, as precision (accepted pairs that are correct) and coverage (correct pairs found). The main models are refitted with deterministic links only, as a sensitivity check.
}

\pt{%
Ficam fora do painel os meses em que o veículo não teve nenhuma viagem e os veículos sem unidade atribuível no mês. Relatos sem veículo identificável não entram no modelo por veículo, mas entram nas taxas por unidade, que podem ser calculadas com todos os relatos; a diferença entre as duas versões é reportada.
}
\en{%
Months in which the vehicle had no trips and vehicles with no assignable site in the month are left out of the panel. Reports without an identifiable vehicle do not enter the vehicle model, but do enter site rates, which can be computed with all reports; the difference between the two versions is reported.
}

\section{\tr{Etapa 2: modelo principal}{Stage 2: main model}}
\label{sec:met-modelo}

\pt{%
O modelo principal é uma binomial negativa hierárquica:
}
\en{%
The main model is a hierarchical negative binomial:
}
\begin{equation}
\label{eq:modelo}
\begin{aligned}
Y_{it} &\sim \text{NB}(\mu_{it}, \phi), \\
\log \mu_{it} &= \log E_{it} + \alpha_{k(i)} + \mathbf{z}_{it}^\top \bm{\beta} + \gamma_{t} + u_{s(i,t)} + v_i, \\
u_s &\sim \mathcal{N}(0, \sigma_u^2), \qquad v_i \sim \mathcal{N}(0, \sigma_v^2),
\end{aligned}
\end{equation}
\pt{%
em que $E_{it}$ são as horas de viagem, $\alpha_{k(i)}$ é o efeito do tipo de uso $k$ do veículo, $\mathbf{z}_{it}$ são as taxas padronizadas de eventos de condução, $\gamma_t$ é um efeito de calendário (efeito fixo de mês ou tendência mais sazonalidade, conforme o número de meses disponíveis), $s(i,t)$ é a unidade do veículo $i$ no mês $t$, e $\phi$ é o parâmetro de dispersão da NB2 (Equação~\ref{eq:nb2}).
}
\en{%
where $E_{it}$ is trip hours, $\alpha_{k(i)}$ is the effect of the vehicle's use type $k$, $\mathbf{z}_{it}$ are the standardized driving event rates, $\gamma_t$ is a calendar effect (month fixed effect or trend plus seasonality, depending on the number of months available), $s(i,t)$ is the site of vehicle $i$ in month $t$, and $\phi$ is the NB2 dispersion parameter (Equation~\ref{eq:nb2}).
}

\pt{%
Três escolhas do modelo vêm dos estudos preliminares (Capítulo~\ref{cap:viabilidade}). O efeito de veículo $v_i$ é obrigatório, porque sem ele os testes sobre telemetria tiveram falsos positivos bem acima do nominal. O tipo de uso $\alpha_k$ é obrigatório, porque o caminhão de manobra tem taxa por hora muito maior que os demais e gera mais eventos de manobra, de modo que sem esse ajuste a telemetria parece associada a acidentes só por acompanhar o tipo de operação. A exposição é o tempo e não a distância, porque, por milha, a comparação entre tipos de uso fica dominada pelo fato de o caminhão de manobra quase não rodar.
}
\en{%
Three model choices come from the preliminary studies (Chapter~\ref{cap:viabilidade}). The vehicle effect $v_i$ is mandatory, because without it the telematics tests had false positives well above the nominal level. The use type $\alpha_k$ is mandatory, because the yard truck has a much higher rate per hour than the others and generates more maneuvering events, so that without this adjustment telematics appears associated with accidents only because it tracks the type of operation. Exposure is time, not distance, because per mile the comparison across use types is dominated by the yard truck barely driving.
}

\pt{%
Os coeficientes de telemetria são lidos como associações, não como efeitos causais: o mesmo motorista, a mesma unidade e o mesmo tipo de manobra podem gerar ao mesmo tempo mais eventos e mais acidentes. Como as taxas de eventos são correlacionadas entre si, cada tipo de evento é testado num modelo próprio (o modelo~\eqref{eq:modelo} com um elemento em $\mathbf{z}_{it}$), e um modelo com todos os tipos juntos é ajustado como sensibilidade.
}
\en{%
Telematics coefficients are read as associations, not causal effects: the same driver, site and type of maneuver can generate both more events and more accidents. Since event rates are correlated with each other, each event type is tested in its own model (model~\eqref{eq:modelo} with one element in $\mathbf{z}_{it}$), and a model with all types together is fitted as a sensitivity check.
}

\pt{%
O ajuste é por máxima verossimilhança com aproximação de Laplace no glmmTMB \parencite{brooks2017}. As unidades e os veículos são cruzados no tempo (um veículo pode mudar de unidade), o que o glmmTMB acomoda com dois efeitos aleatórios não aninhados. A hipótese de proporcionalidade do \textit{offset} é verificada substituindo-o por $\theta \log E_{it}$ com $\theta$ livre (Seção~\ref{sec:fund-contagem}).
}
\en{%
Fitting is by maximum likelihood with the Laplace approximation in glmmTMB \parencite{brooks2017}. Sites and vehicles are crossed over time (a vehicle may move between sites), which glmmTMB accommodates with two non-nested random effects. The proportionality assumption of the offset is checked by replacing it with $\theta \log E_{it}$ with free $\theta$ (Section~\ref{sec:fund-contagem}).
}

\section{\tr{Modelos comparados e diagnóstico}{Compared models and diagnostics}}
\label{sec:met-comparados}

\pt{%
O objetivo (b) compara o modelo principal a alternativas mais simples e mais complexas (Tabela~\ref{tab:modelos}). A comparação é feita pelo AIC no ajuste e pelo escore logarítmico na validação temporal (Seção~\ref{sec:met-validacao}). As variâncias dos efeitos aleatórios são testadas por razão de verossimilhanças; como $\sigma^2 = 0$ está na fronteira do espaço de parâmetros, a distribuição de referência é uma mistura meio a meio de $\chi^2_0$ e $\chi^2_1$, o que equivale a dividir por dois o valor-$p$ usual.
}
\en{%
Objective (b) compares the main model with simpler and more complex alternatives (Table~\ref{tab:modelos}). The comparison uses AIC in fitting and the log score in temporal validation (Section~\ref{sec:met-validacao}). Random effect variances are tested by likelihood ratio; since $\sigma^2 = 0$ lies on the boundary of the parameter space, the reference distribution is an equal mixture of $\chi^2_0$ and $\chi^2_1$, which amounts to halving the usual $p$-value.
}

\begin{table}[ht]
\centering\small
\caption{\tr{Modelos comparados. Todos têm as horas de viagem como \textit{offset}.}{Compared models. All have trip hours as offset.}}
\label{tab:modelos}
\begin{tabularx}{\textwidth}{l l >{\raggedright\arraybackslash}X}
\toprule
& \tr{Distribuição}{Distribution} & \tr{Preditor}{Predictor} \\
\midrule
M0 & Poisson & \tr{tipo de uso, calendário}{use type, calendar} \\
M1 & NB2 & \tr{como M0}{as M0} \\
M2 & NB2 & \tr{M1 + efeito de veículo $v_i$}{M1 + vehicle effect $v_i$} \\
M3 & NB2 & \tr{M2 + efeito de unidade $u_s$ (modelo de referência, sem telemetria)}{M2 + site effect $u_s$ (reference model, without telematics)} \\
M4 & NB2 & \tr{M3 + eventos de condução (modelo~\eqref{eq:modelo})}{M3 + driving events (model~\eqref{eq:modelo})} \\
M5 & \tr{NB2 com inflação de zeros}{Zero-inflated NB2} & \tr{como M3, com probabilidade de zero estrutural constante}{as M3, with constant structural zero probability} \\
M6 & NB2 & \tr{M3 com distância no lugar das horas como exposição}{M3 with distance instead of hours as exposure} \\
\bottomrule
\end{tabularx}
\end{table}

\pt{%
O diagnóstico usa os resíduos quantílicos simulados do DHARMa \parencite{dunn1996,hartig2026}: uniformidade (teste de Kolmogorov e Smirnov), dispersão (razão entre a variância observada e a simulada), excesso de zeros (zeros observados contra simulados) e autocorrelação temporal dentro do veículo. Na exploração preliminar (Seção~\ref{sec:viab-real}), M2 não mostrou superdispersão residual nem excesso de zeros, o que já indica que M5 não deve ser necessário.
}
\en{%
Diagnostics use DHARMa's simulated quantile residuals \parencite{dunn1996,hartig2026}: uniformity (Kolmogorov-Smirnov test), dispersion (ratio of observed to simulated variance), excess zeros (observed against simulated zeros) and within-vehicle temporal autocorrelation. In the preliminary exploration (Section~\ref{sec:viab-real}), M2 showed no residual overdispersion or excess zeros, which already suggests that M5 should not be needed.
}

\section{\tr{Indicadores de risco por unidade}{Site risk indicators}}
\label{sec:met-indicadores}

\pt{%
Para o objetivo (e) e para H1, cada unidade $s$ recebe três indicadores num período: a contagem bruta $Y_s$; a taxa bruta $Y_s/E_s$; e a taxa ajustada, que é a taxa de Bayes empírico da Equação~\eqref{eq:eb} na versão simples ou, na versão completa, a razão de taxas da unidade $\exp(\hat{u}_s)$ estimada pelo modelo M3, com intervalo de previsão do efeito aleatório. A versão completa ajusta também pelo tipo de uso dos veículos da unidade, o que a versão simples não faz.
}
\en{%
For objective (e) and H1, each site $s$ receives three indicators in a period: the raw count $Y_s$; the raw rate $Y_s/E_s$; and the adjusted rate, which is the empirical Bayes rate of Equation~\eqref{eq:eb} in the simple version or, in the full version, the site rate ratio $\exp(\hat{u}_s)$ estimated by model M3, with a prediction interval for the random effect. The full version also adjusts for the use type of the site's vehicles, which the simple one does not.
}

\pt{%
O gráfico de funil \parencite{spiegelhalter2005} plota a razão entre observado e esperado de cada unidade, $Y_s/e_s$, contra o esperado $e_s$, que é a soma, sobre os veículos e meses da unidade, dos valores ajustados pelo modelo sem efeito de unidade. Os limites de 95\% e 99,8\% vêm dos quantis da Poisson com média $e_s$, divididos por $e_s$. As unidades são classificadas em três grupos: acima do esperado (intervalo de $\exp(\hat{u}_s)$ todo acima de 1), compatível com a média e abaixo do esperado. Para mostrar a incerteza do ranking, as posições são recalculadas em 2.000 simulações da distribuição dos efeitos de unidade, e cada unidade recebe um intervalo de posição \parencite{goldstein1996}.
}
\en{%
The funnel plot \parencite{spiegelhalter2005} plots each site's observed-to-expected ratio, $Y_s/e_s$, against the expected count $e_s$, which is the sum over the site's vehicles and months of the fitted values from the model without a site effect. The 95\% and 99.8\% limits come from Poisson quantiles with mean $e_s$, divided by $e_s$. Sites are classified into three groups: above expected (interval of $\exp(\hat{u}_s)$ entirely above 1), compatible with the mean, and below expected. To show ranking uncertainty, positions are recomputed in 2,000 simulations from the distribution of site effects, and each site gets a position interval \parencite{goldstein1996}.
}

\section{\tr{Etapa 3: validação fora da amostra}{Stage 3: out-of-sample validation}}
\label{sec:met-validacao}

\pt{%
A validação tem duas formas. Na \textbf{temporal}, os modelos são ajustados nos meses iniciais e avaliados nos últimos meses, com origem móvel \parencite{tashman2000}: a cada corte, o modelo é reajustado e prevê o mês seguinte. Na \textbf{por unidade}, cada unidade é deixada de fora por vez, e seus veículos são previstos com o efeito de unidade igual a zero; ela mede quanto o modelo transporta para uma unidade nova, o caso de um contrato que começa.
}
\en{%
Validation takes two forms. In the \textbf{temporal} one, models are fitted on the early months and evaluated on the last months, with a rolling origin \parencite{tashman2000}: at each cutoff, the model is refitted and forecasts the next month. In the \textbf{site} one, each site is left out in turn, and its vehicles are predicted with the site effect set to zero; it measures how well the model transfers to a new site, the case of a contract that starts.
}

\pt{%
As métricas são: (i) o escore logarítmico médio por veículo e mês, $\overline{S} = n^{-1}\sum \log p(y_{it} \mid \hat{\mu}_{it}, \hat{\phi})$, em que a previsão usa os efeitos de veículo e unidade estimados no treino; (ii) a calibração, pela transformada integral de probabilidade não aleatorizada e pelo diagrama de calibração marginal \parencite{czado2009}; (iii) a concentração, a fração dos acidentes do período de teste que ocorreu nos 10\% e 20\% de veículos com maior taxa prevista, uma métrica que a gestão entende diretamente.
}
\en{%
The metrics are: (i) the mean log score per vehicle-month, $\overline{S} = n^{-1}\sum \log p(y_{it} \mid \hat{\mu}_{it}, \hat{\phi})$, where the forecast uses the vehicle and site effects estimated in training; (ii) calibration, through the nonrandomized probability integral transform and the marginal calibration diagram \parencite{czado2009}; (iii) concentration, the share of the test period's accidents that occurred in the 10\% and 20\% of vehicles with the highest predicted rate, a metric management understands directly.
}

\section{\tr{Testes das hipóteses}{Hypothesis tests}}
\label{sec:testes}

\pt{%
Todos os testes usam $\alpha = 0{,}05$. Os intervalos por bootstrap usam 2.000 reamostragens, com o grupo reamostrado indicado em cada caso, para respeitar a dependência dos dados.
}
\en{%
All tests use $\alpha = 0.05$. Bootstrap intervals use 2,000 resamples, with the resampled group indicated in each case, to respect the dependence in the data.
}

\pt{%
\begin{description}[leftmargin=0.6cm,style=nextline]
  \item[H1a] Os indicadores calculados num semestre preveem a taxa observada no semestre seguinte. O erro de cada indicador é $\sum_s w_s \, |\hat{r}_s - r_s^{+}|$, com $r_s^{+}$ a taxa do semestre seguinte e pesos $w_s$ proporcionais à exposição desse semestre. A estatística é a diferença de erro entre a taxa bruta e a ajustada, com intervalo e valor-$p$ unilateral por bootstrap pareado por unidade. A contagem bruta não entra nesse teste, porque está em outra escala.
  \item[H1b] Teste de homogeneidade de Poisson, $X^2 = \sum_s (Y_s - e_s)^2/e_s$, comparado a $\chi^2_{S-1}$, e teste de razão de verossimilhanças de $\sigma_u^2 = 0$ em M3. Como complemento, a fração de unidades fora do limite de 95\% do funil é comparada a 5\% por um teste binomial.
  \item[H1c] Correlação de postos de Spearman entre cada indicador e a exposição da unidade, com intervalo de 90\% por bootstrap por unidade. A equivalência é declarada se o intervalo couber em $(-0{,}3;\ 0{,}3)$, o que equivale aos dois testes unilaterais a 5\%.
  \item[H2a] Teste de Wald de $\beta_j = 0$ em cada um dos modelos por tipo de evento, com correção de \textcite{holm1979} sobre os tipos. Na parte preditiva, a diferença de escore logarítmico médio entre M4 e M3 na validação temporal, com intervalo e valor-$p$ por bootstrap por veículo, no espírito de \textcite{diebold1995}.
  \item[H2b] Para cada tipo de evento, o intervalo de 90\% da razão de taxas $\exp(\beta_j)$ é comparado à margem $[0{,}8;\ 1{,}25]$. Cada tipo é classificado como associado (rejeita H2a após Holm), desprezível (intervalo dentro da margem) ou inconclusivo.
  \item[H3a] A probabilidade de relato $p$ é estimada pelo modelo logístico da Seção~\ref{sec:met-subnotificacao}. Como erros de ligação também produzem ordens sem relato, $H_0$ é rejeitada se o limite superior do intervalo de 95\% de $p$ ficar abaixo de $1 - \varepsilon$, em que $\varepsilon$ é a taxa de pares perdidos pela ligação, medida na amostra revisada.
  \item[H3b] Teste de razão de verossimilhanças da variância do efeito de unidade no modelo logístico, com a correção de fronteira da Seção~\ref{sec:met-comparados}.
\end{description}
}
\en{%
\begin{description}[leftmargin=0.6cm,style=nextline]
  \item[H1a] The indicators computed in one semester predict the rate observed in the next. Each indicator's error is $\sum_s w_s \, |\hat{r}_s - r_s^{+}|$, with $r_s^{+}$ the next semester's rate and weights $w_s$ proportional to that semester's exposure. The statistic is the error difference between the raw and adjusted rates, with interval and one-sided $p$-value from a bootstrap paired by site. The raw count does not enter this test because it is on a different scale.
  \item[H1b] Poisson homogeneity test, $X^2 = \sum_s (Y_s - e_s)^2/e_s$, compared with $\chi^2_{S-1}$, and likelihood ratio test of $\sigma_u^2 = 0$ in M3. As a complement, the share of sites outside the funnel's 95\% limit is compared with 5\% by a binomial test.
  \item[H1c] Spearman rank correlation between each indicator and site exposure, with a 90\% bootstrap interval by site. Equivalence is declared if the interval fits within $(-0.3,\ 0.3)$, which is equivalent to the two one-sided tests at 5\%.
  \item[H2a] Wald test of $\beta_j = 0$ in each event-type model, with \textcite{holm1979} correction across types. In the predictive part, the difference in mean log score between M4 and M3 in temporal validation, with interval and $p$-value from a bootstrap by vehicle, in the spirit of \textcite{diebold1995}.
  \item[H2b] For each event type, the 90\% interval of the rate ratio $\exp(\beta_j)$ is compared with the margin $[0.8,\ 1.25]$. Each type is classified as associated (rejects H2a after Holm), negligible (interval within the margin) or inconclusive.
  \item[H3a] The reporting probability $p$ is estimated by the logistic model of Section~\ref{sec:met-subnotificacao}. Since linkage errors also produce orders without reports, $H_0$ is rejected if the upper limit of the 95\% interval of $p$ is below $1 - \varepsilon$, where $\varepsilon$ is the rate of pairs missed by linkage, measured on the reviewed sample.
  \item[H3b] Likelihood ratio test of the site effect variance in the logistic model, with the boundary correction of Section~\ref{sec:met-comparados}.
\end{description}
}

\pt{%
As margens de equivalência seguem convenções usuais: a razão de taxas entre 0,8 e 1,25 é simétrica na escala log e é a margem habitual em estudos de bioequivalência, e uma correlação abaixo de 0,3 em módulo costuma ser descrita como fraca \preencher{citar a fonte de cada convenção}. Elas serão revistas com a gestão de segurança, para refletir o menor efeito que mudaria uma decisão, e registradas antes da análise final \preencher{registrar margens e plano de análise, por exemplo no OSF, antes de ver os dados finais}.
}
\en{%
The equivalence margins follow common conventions: a rate ratio between 0.8 and 1.25 is symmetric on the log scale and is the usual margin in bioequivalence studies, and a correlation below 0.3 in absolute value is usually described as weak \preencher{cite the source of each convention}. They will be reviewed with safety management, to reflect the smallest effect that would change a decision, and registered before the final analysis \preencher{register margins and analysis plan, e.g.\ on OSF, before seeing the final data}.
}

\section{\tr{Etapa 4: subnotificação}{Stage 4: underreporting}}
\label{sec:met-subnotificacao}

\pt{%
\textbf{Detecção automática de colisão.} A primeira parte do objetivo (d) mede a cobertura da detecção automática da telemetria como fonte de acidentes de pátio: a fração dos acidentes relatados com ao menos uma detecção no mesmo veículo numa janela de dois dias, por tipo de uso, com intervalo de Wilson, e a fração das detecções descartadas como falso alarme na revisão. A exploração preliminar indicou cobertura muito baixa (Seção~\ref{sec:viab-real}), e por isso essa fonte não é usada na captura e recaptura.
}
\en{%
\textbf{Automatic crash detection.} The first part of objective (d) measures the coverage of telematics automatic detection as a source of yard accidents: the share of reported accidents with at least one detection on the same vehicle within a two-day window, by use type, with a Wilson interval, and the share of detections dismissed as false alarms in review. The preliminary exploration indicated very low coverage (Section~\ref{sec:viab-real}), so this source is not used for capture-recapture.
}

\pt{%
\textbf{Manutenção como segunda fonte.} Sejam A os relatos de colisão e B as ordens de manutenção abertas por dano ou acidente. Uma ordem $j$ do veículo $i$ na data $d_j$ é pareada a um relato do mesmo veículo com data em $[d_j - \Delta_1,\ d_j + \Delta_2]$; a janela principal e as alternativas são fixadas a partir da distribuição dos intervalos entre relato e ordem nos pares inequívocos \preencher{definir $\Delta_1$ e $\Delta_2$}. Para cada ordem, $R_j = 1$ se há relato pareado. A probabilidade de relato é modelada por
}
\en{%
\textbf{Maintenance as second source.} Let A be the collision reports and B the maintenance orders opened for damage or accident. An order $j$ of vehicle $i$ on date $d_j$ is matched to a report for the same vehicle dated within $[d_j - \Delta_1,\ d_j + \Delta_2]$; the main and alternative windows are set from the distribution of report-to-order intervals in unambiguous pairs \preencher{define $\Delta_1$ and $\Delta_2$}. For each order, $R_j = 1$ if there is a matched report. The reporting probability is modeled by
}
\begin{equation}
\label{eq:relato}
\text{logit}\, \Pr(R_j = 1) = \eta + \mathbf{w}_j^\top \bm{\delta} + a_{s(j)}, \qquad a_s \sim \mathcal{N}(0, \sigma_a^2),
\end{equation}
\pt{%
em que $\mathbf{w}_j$ reúne o tipo de uso, o custo ou a classe do reparo e o período, e $a_s$ é o efeito de unidade. A probabilidade de relato da unidade, $p_s$, é estimada com encolhimento, porque a estimativa bruta por unidade se mostrou muito ruidosa no estudo de simulação (Seção~\ref{sec:poder}).
}
\en{%
where $\mathbf{w}_j$ holds use type, repair cost or class and period, and $a_s$ is the site effect. The site reporting probability $p_s$ is estimated with shrinkage, because the raw site estimate proved very noisy in the simulation study (Section~\ref{sec:poder}).
}

\pt{%
\textbf{Interpretação como limite.} Nem todo dano gera um acidente relatável (desgaste, dano preexistente, vandalismo), e nem todo acidente gera uma ordem (dano só na carreta do cliente, na doca ou na estrutura). Assim, $p_s$ é a probabilidade de relato \emph{entre os acidentes que danificaram o próprio caminhão}. O total de acidentes é estimado pela Equação~\eqref{eq:lp}, com $n_1$ relatos, $n_2$ ordens e $m$ pares, estratificando por tipo de acidente e gravidade para reduzir a heterogeneidade. Como acidentes mais graves tendem a ser relatados e reparados ao mesmo tempo, a dependência entre as fontes deve ser positiva, e o total estimado deve ser lido como limite inferior \parencite{hook1995}; a sensibilidade a uma dependência de magnitude fixada é mostrada explicitamente.
}
\en{%
\textbf{Interpretation as a bound.} Not every damage stems from a reportable accident (wear, pre-existing damage, vandalism), and not every accident generates an order (damage only to the client's trailer, the dock or a structure). Hence $p_s$ is the reporting probability \emph{among accidents that damaged the truck itself}. The total number of accidents is estimated by Equation~\eqref{eq:lp}, with $n_1$ reports, $n_2$ orders and $m$ pairs, stratifying by accident type and severity to reduce heterogeneity. Since more severe accidents tend to be both reported and repaired, dependence between sources should be positive, and the estimated total should be read as a lower bound \parencite{hook1995}; sensitivity to a fixed magnitude of dependence is shown explicitly.
}

\section{\tr{Plano alternativo}{Fallback plan}}
\label{sec:met-alternativo}

\pt{%
Se o uso da telemetria for restrito, o mesmo arcabouço é aplicado por unidade e mês, com as horas trabalhadas da folha como exposição: $\log \mu_{st} = \log H_{st} + \gamma_t + \mathbf{x}_{st}^\top \bm{\beta} + u_s$, com $\mathbf{x}_{st}$ podendo incluir a fração de horas extras. O objetivo (b), o objetivo (e) e a hipótese H1 seguem testáveis. O objetivo (c) e a hipótese H2 deixam de existir. A parte de manutenção do objetivo (d) e a hipótese H3 continuam possíveis se o cadastro da manutenção associar cada veículo a uma unidade; a medida de cobertura da detecção automática deixa de existir \preencher{confirmar se o cadastro da manutenção traz a unidade do veículo}.
}
\en{%
If the use of telematics is restricted, the same framework is applied per site and month, with payroll hours worked as exposure: $\log \mu_{st} = \log H_{st} + \gamma_t + \mathbf{x}_{st}^\top \bm{\beta} + u_s$, where $\mathbf{x}_{st}$ may include the share of overtime hours. Objective (b), objective (e) and hypothesis H1 remain testable. Objective (c) and hypothesis H2 are dropped. The maintenance part of objective (d) and hypothesis H3 remain possible if the maintenance registry assigns each vehicle to a site; the coverage measure of automatic detection is dropped \preencher{confirm whether the maintenance registry carries the vehicle's site}.
}

\section{\tr{Análises de sensibilidade}{Sensitivity analyses}}
\label{sec:met-sensibilidade}

\pt{%
\begin{itemize}
  \item \textbf{Desfecho:} só colisões com lesão; colisões mais quase-acidentes; só relatos com revisão concluída.
  \item \textbf{Exposição:} distância no lugar das horas (M6); horas da folha no lugar das horas de viagem, por unidade e mês.
  \item \textbf{Ligação:} só ligações determinísticas; janelas alternativas na subnotificação.
  \item \textbf{Período:} exclusão dos meses de implantação do aplicativo de relatos.
  \item \textbf{Covariáveis de organização do trabalho:} fração de horas extras e, se disponível, fração de recém-contratados por unidade e mês \parencite{dembe2005,breslin2006}.
  \item \textbf{Telemetria:} todos os tipos de evento num único modelo; erros-padrão robustos por veículo \parencite{cameron2015} no lugar do efeito aleatório.
\end{itemize}
}
\en{%
\begin{itemize}
  \item \textbf{Outcome:} injury collisions only; collisions plus near misses; only reports with completed review.
  \item \textbf{Exposure:} distance instead of hours (M6); payroll hours instead of trip hours, per site and month.
  \item \textbf{Linkage:} deterministic links only; alternative windows in underreporting.
  \item \textbf{Period:} excluding the reporting application's roll-out months.
  \item \textbf{Work organization covariates:} share of overtime hours and, if available, share of new hires per site and month \parencite{dembe2005,breslin2006}.
  \item \textbf{Telematics:} all event types in a single model; vehicle-clustered robust standard errors \parencite{cameron2015} instead of the random effect.
\end{itemize}
}

\section{\tr{Extensão opcional: previsão da exposição}{Optional extension: exposure forecasting}}
\label{sec:met-previsao}

\pt{%
As horas de operação de cada unidade são previstas como distribuição, semanal e mensalmente, por três métodos: ingênuo sazonal, alisamento exponencial com tendência amortecida \parencite{gardner1985} e um modelo de séries hierárquicas, com reconciliação entre unidades, regiões e total \parencite{hyndman2021}. A avaliação usa origem móvel, CRPS, perda quantílica e cobertura dos intervalos de 80\% e 95\% \parencite{gneiting2007}.
}
\en{%
Each site's operating hours are forecast as a distribution, weekly and monthly, by three methods: seasonal naive, exponential smoothing with damped trend \parencite{gardner1985} and a hierarchical time series model, with reconciliation across sites, regions and total \parencite{hyndman2021}. Evaluation uses a rolling origin, CRPS, quantile loss and coverage of the 80\% and 95\% intervals \parencite{gneiting2007}.
}

\pt{%
Para projetar acidentes, a previsão da exposição $\tilde{E}$ é combinada à taxa estimada $\hat{\lambda}$: $Y \mid \tilde{E} \sim \text{NB}(\hat{\lambda}\tilde{E}, \hat{\phi})$. Pela lei da variância total, $\text{Var}(Y) = \mathbb{E}[\text{Var}(Y \mid \tilde{E})] + \hat{\lambda}^2 \, \text{Var}(\tilde{E})$, e a fração da variância devida à exposição é a segunda parcela dividida pelo total. No estudo preliminar, essa fração foi de cerca de 1\% para uma unidade com um acidente esperado por mês (Seção~\ref{sec:previsao}); ela cresce com o número esperado de acidentes, e por isso a projeção é feita só em níveis agregados (região e trimestre).
}
\en{%
To project accidents, the exposure forecast $\tilde{E}$ is combined with the estimated rate $\hat{\lambda}$: $Y \mid \tilde{E} \sim \text{NB}(\hat{\lambda}\tilde{E}, \hat{\phi})$. By the law of total variance, $\text{Var}(Y) = \mathbb{E}[\text{Var}(Y \mid \tilde{E})] + \hat{\lambda}^2 \, \text{Var}(\tilde{E})$, and the share of variance due to exposure is the second term divided by the total. In the preliminary study, this share was about 1\% for a site with one expected accident per month (Section~\ref{sec:previsao}); it grows with the expected number of accidents, which is why projection is done only at aggregate levels (region and quarter).
}

\section{\tr{Ferramentas e reprodutibilidade}{Tools and reproducibility}}
\label{sec:met-ferramentas}

\pt{%
A extração e a agregação rodam dentro do ambiente de dados da empresa, em SQL, e só tabelas agregadas por veículo ou unidade e mês, sem dados pessoais, saem para a análise. A modelagem usa Python (preparação, bootstrap, previsão e figuras) e R (modelos mistos com glmmTMB e diagnóstico com DHARMa, num ambiente renv com versões fixas), com código versionado e sementes fixas. Toda tabela e figura do texto é gerada por script, e os números citados no texto são conferidos contra as tabelas geradas. O código e o estudo de simulação podem ser publicados; os dados da operação, não.
}
\en{%
Extraction and aggregation run inside the company's data environment, in SQL, and only aggregated tables per vehicle or site and month, with no personal data, leave for analysis. Modeling uses Python (preparation, bootstrap, forecasting and figures) and R (mixed models with glmmTMB and diagnostics with DHARMa, in an renv environment with pinned versions), with versioned code and fixed seeds. Every table and figure in the text is generated by script, and the numbers cited in the text are checked against the generated tables. The code and the simulation study can be published; the operational data cannot.
}

\chapter{\tr{Estudos preliminares de viabilidade}{Preliminary feasibility studies}}
\label{cap:viabilidade}

\pt{%
Antes de ter acesso aos dados da operação, dois estudos usam a estrutura de um conjunto público (Capítulo~\ref{cap:dados}) para responder perguntas de viabilidade: (i) com o volume de acidentes esperado, os modelos têm poder para testar as hipóteses? (ii) a previsão da exposição, base do objetivo opcional (f), é útil? Os números deste capítulo vêm de \texttt{analise/viabilidade/poder.py} e \texttt{previsao.py} e são de dados simulados: indicam ordens de grandeza, não resultados. A Seção~\ref{sec:viab-real} completa o capítulo com uma exploração agregada dos dados da operação, que verifica se os dados reais sustentam as hipóteses.
}
\en{%
Before accessing operational data, two studies use the structure of a public dataset (Chapter~\ref{cap:dados}) to answer feasibility questions: (i) given the expected accident volume, do the models have power to test the hypotheses? (ii) is the exposure forecast, which underlies optional objective (f), useful? Numbers in this chapter come from \texttt{analise/viabilidade/poder.py} and \texttt{previsao.py} and are from simulated data: they indicate orders of magnitude, not results. Section~\ref{sec:viab-real} completes the chapter with an aggregated exploration of the operational data, which checks whether the real data support the hypotheses.
}

\section{\tr{Poder estatístico (semi-sintético)}{Statistical power (semi-synthetic)}}
\label{sec:poder}

\pt{%
A exposição (milhas por caminhão e mês) e a estrutura de caminhões e terminais vêm do conjunto público; os acidentes são simulados com efeitos conhecidos: efeito de telemetria $\beta \in \{0;\ 0{,}15;\ 0{,}30\}$ no log da taxa por desvio padrão, efeito aleatório de site (desvio padrão 0,40) e de caminhão (0,30), superdispersão adicional, detecção de 35\% das colisões pelos sensores e probabilidade de relato por site com média 0,65. O total esperado de acidentes relatados no período foi fixado em 70, 300 e 600. Duas bases: o painel do conjunto como está (92 caminhões, 24 terminais, 36 meses) e uma base na \emph{escala do projeto} (280 caminhões em 40 sites por 14 meses, com milhas reamostradas). O desenho segue a estrutura usual de estudos de simulação: objetivos, mecanismo gerador dos dados, estimandos, métodos e medidas de desempenho \parencite{morris2019}. Foram 200 repetições por cenário; o modelo ajustado é a binomial negativa com efeito fixo de site e milhas como \textit{offset}, ainda sem efeito aleatório de caminhão.
}
\en{%
Exposure (miles per truck and month) and the truck and terminal structure come from the public dataset; accidents are simulated with known effects: telematics effect $\beta \in \{0, 0.15, 0.30\}$ on the log rate per standard deviation, site (standard deviation 0.40) and truck (0.30) random effects, extra overdispersion, sensor detection of 35\% of collisions and a site reporting probability with mean 0.65. The expected number of reported accidents in the period was set to 70, 300 and 600. Two bases: the dataset panel as is (92 trucks, 24 terminals, 36 months) and a base at \emph{project scale} (280 trucks in 40 sites over 14 months, with resampled miles). The design follows the usual structure of simulation studies: aims, data-generating mechanism, estimands, methods and performance measures \parencite{morris2019}. There were 200 replications per scenario; the fitted model is the negative binomial with site fixed effects and miles as offset, still without a truck random effect.
}

\begin{table}[ht]
\centering\small
\caption{\tr{Proporção de repetições com $p < 0{,}05$ para o efeito de telemetria. Com $\beta = 0$, é a taxa de falso positivo (o ideal é 0,05).}%
            {Share of replications with $p < 0.05$ for the telematics effect. With $\beta = 0$, it is the false positive rate (ideal: 0.05).}}
\label{tab:poder}
\begin{tabular}{l r r r r}
\toprule
\tr{Base}{Base} & \tr{Acidentes}{Accidents} & $\beta = 0$ & $\beta = 0{,}15$ & $\beta = 0{,}30$ \\
\midrule
\tabinput{tabelas/tab_poder.tex}
\bottomrule
\end{tabular}
\end{table}

\begin{figure}[ht]
\centering
\includegraphics[width=\textwidth]{\arqfig{fig_poder}}
\caption{\tr{Poder para detectar o efeito de telemetria. Linhas cinza: 0,80 e 0,05.}{Power to detect the telematics effect. Grey lines: 0.80 and 0.05.}}
\label{fig:poder}
\end{figure}

\pt{%
\textbf{H2 (telemetria).} Na escala do projeto, com cerca de 300 acidentes relatados, o poder é 0,75 para um efeito moderado ($\beta = 0{,}15$, cerca de 16\% a mais de risco por desvio padrão) e praticamente 1 para $\beta = 0{,}30$; com 600 acidentes, 0,89 e 1. Com 70 acidentes, só um efeito grande tem chance razoável (0,53). \textbf{Implicação:} o desfecho principal precisa incluir todas as colisões com dano relatadas, não só os acidentes graves; uma definição estreita derruba o poder.
}
\en{%
\textbf{H2 (telematics).} At project scale, with about 300 reported accidents, power is 0.75 for a moderate effect ($\beta = 0.15$, about 16\% more risk per standard deviation) and practically 1 for $\beta = 0.30$; with 600 accidents, 0.89 and 1. With 70 accidents, only a large effect has a reasonable chance (0.53). \textbf{Implication:} the main outcome must include all reported damage collisions, not only severe accidents; a narrow definition kills power.
}

\pt{%
\textbf{Falso positivo.} Na base do Kaggle com 600 acidentes, o efeito nulo saiu significativo em 15\% das repetições, três vezes o nominal. A causa provável é ignorar a correlação entre os meses do mesmo caminhão (36 meses por caminhão nessa base), justamente o que o efeito aleatório de veículo do modelo hierárquico trata. \textbf{Implicação:} o modelo final precisa do efeito aleatório de caminhão ou de erros-padrão robustos por caminhão; a versão sem ele não serve para inferência.
}
\en{%
\textbf{False positives.} In the Kaggle base with 600 accidents, the null effect was significant in 15\% of replications, three times the nominal rate. The likely cause is ignoring the correlation between months of the same truck (36 months per truck in this base), exactly what the hierarchical model's vehicle random effect handles. \textbf{Implication:} the final model needs a truck random effect or truck-clustered robust standard errors; the version without it is not fit for inference.
}

\begin{table}[ht]
\centering\small
\caption{\tr{H1: correlação de postos (Spearman) entre o indicador por site e o risco verdadeiro do site; média das repetições.}%
            {H1: rank correlation (Spearman) between the site indicator and the site's true risk; mean over replications.}}
\label{tab:poder_h1}
\begin{tabular}{l r r r r}
\toprule
\tr{Base}{Base} & \tr{Acidentes}{Accidents} & \tr{Contagem}{Count} & \tr{Taxa bruta}{Raw rate} & \tr{Bayes empírico}{Empirical Bayes} \\
\midrule
\tabinput{tabelas/tab_poder_h1.tex}
\bottomrule
\end{tabular}
\end{table}

\pt{%
\textbf{H1 (ranking de sites).} A contagem bruta é o pior indicador (correlação de 0,30 com o risco verdadeiro na escala do projeto com 300 acidentes), porque mede sobretudo o tamanho do site. Taxa bruta e Bayes empírico ficam bem acima (0,51 e 0,54), com vantagem pequena, mas consistente, do Bayes empírico na escala do projeto. Mesmo com 600 acidentes, a correlação não passa de cerca de 0,6. \textbf{Implicação:} ranking de 40 sites com um ano de dados é ruidoso por natureza; a dissertação deve apresentar intervalos (gráfico de funil) e grupos de risco, não posições exatas, e a contribuição de H1 está mais em trocar a contagem pela taxa ajustada do que no ganho do encolhimento.
}
\en{%
\textbf{H1 (site ranking).} The raw count is the worst indicator (correlation of 0.30 with true risk at project scale with 300 accidents), because it mostly measures site size. Raw rate and empirical Bayes are well above (0.51 and 0.54), with a small but consistent advantage for empirical Bayes at project scale. Even with 600 accidents, the correlation does not exceed about 0.6. \textbf{Implication:} ranking 40 sites with one year of data is inherently noisy; the dissertation should present intervals (funnel plot) and risk groups, not exact positions, and H1's contribution lies more in replacing counts with adjusted rates than in the gain from shrinkage.
}

\begin{table}[ht]
\centering\small
\caption{\tr{H3: erro absoluto médio da probabilidade de relato por site.}{H3: mean absolute error of the site reporting probability.}}
\label{tab:poder_h3}
\begin{tabular}{l r r r r}
\toprule
\tr{Base}{Base} & \tr{Acidentes}{Accidents} & \tr{Colisões detectadas}{Detected collisions} & \tr{Bruta}{Raw} & \tr{Com encolhimento}{Shrunk} \\
\midrule
\tabinput{tabelas/tab_poder_h3.tex}
\bottomrule
\end{tabular}
\end{table}

\pt{%
\textbf{H3 (subnotificação).} A estimativa bruta da probabilidade de relato por site erra em média 0,23 na escala do projeto com 300 acidentes (cerca de 160 colisões detectadas); o encolhimento reduz o erro para 0,11. \textbf{Implicação:} estimar a subnotificação por site é viável só com encolhimento, e a média geral de relato é bem mais precisa que os valores por site.
}
\en{%
\textbf{H3 (underreporting).} The raw site reporting probability estimate errs by 0.23 on average at project scale with 300 accidents (about 160 detected collisions); shrinkage reduces the error to 0.11. \textbf{Implication:} estimating underreporting per site is feasible only with shrinkage, and the overall reporting rate is much more precise than site values.
}

\section{\tr{Previsão da exposição (objetivo f)}{Exposure forecasting (objective f)}}
\label{sec:previsao}

\pt{%
Com as horas de viagem semanais de \nsitesprev{} cidades de origem do conjunto público (\nsemanas{} semanas), as horas das quatro semanas seguintes foram previstas como distribuição, em origem móvel ao longo do último ano, por três métodos.
}
\en{%
Using weekly trip hours from \nsitesprev{} origin cities in the public dataset (\nsemanas{} weeks), hours for the next four weeks were forecast as a distribution, with a rolling origin over the last year, by three methods.
}

\begin{table}[ht]
\centering\small
\caption{\tr{Previsão das horas das 4 semanas seguintes: perda quantílica média relativa ($q = 0{,}1;\ 0{,}5;\ 0{,}9$), cobertura do intervalo de 80\% e erro relativo mediano da mediana.}%
            {Forecast of the next 4 weeks' hours: mean relative quantile loss ($q = 0.1, 0.5, 0.9$), 80\% interval coverage and median relative error of the median.}}
\label{tab:previsao}
\begin{tabular}{l r r r}
\toprule
\tr{Método}{Method} & \tr{Perda quantílica}{Quantile loss} & \tr{Cobertura 80\%}{80\% coverage} & \tr{Erro da mediana}{Median error} \\
\midrule
\tabinput{tabelas/tab_previsao.tex}
\bottomrule
\end{tabular}
\end{table}

\begin{figure}[ht]
\centering
\includegraphics[width=\textwidth]{\arqfig{fig_previsao}}
\caption{\tr{Previsão das horas do mês seguinte no site de maior volume (dados públicos).}{Next month's hours forecast at the highest-volume site (public data).}}
\label{fig:previsao}
\end{figure}

\pt{%
O ETS amortecido foi o melhor nas três métricas, com erro mediano de cerca de 8\% e cobertura um pouco abaixo da nominal. O achado mais importante, porém, é outro: para um site com cerca de 1 acidente esperado por mês (a ordem de grandeza do projeto), só cerca de \fracexposicao\% da variância do número de acidentes vem da incerteza sobre a exposição; o resto é a variação de Poisson. \textbf{Implicação:} prever a exposição quase não melhora a projeção de acidentes de um site no mês seguinte. O objetivo (f) deve ser reposicionado: usar a previsão de exposição para planejamento de escala e supervisão em si, ou projetar acidentes em horizontes e agregados maiores (trimestre, região), onde a exposição pesa mais. Decisão (2026-11-08): a extensão passa a ser a previsão da exposição para planejamento de escala e supervisão, com projeção de acidentes só em níveis agregados; o pré-projeto já reflete isso.
}
\en{%
Damped ETS was best on all three metrics, with a median error of about 8\% and coverage slightly below nominal. The most important finding, however, is another: for a site with about 1 expected accident per month (the project's order of magnitude), only about \fracexposicao\% of the variance of the accident count comes from exposure uncertainty; the rest is Poisson variation. \textbf{Implication:} forecasting exposure barely improves next month's accident projection for a site. Objective (f) should be repositioned: use the exposure forecast for staffing and supervision planning in itself, or project accidents over longer horizons and larger aggregates (quarter, region), where exposure matters more. Decision (2026-11-08): the extension becomes the exposure forecast for staffing and supervision planning, with accident projection only at aggregate levels; the pre-proposal already reflects this.
}

\section{\tr{Exploração dos dados da operação}{Exploration of operational data}}
\label{sec:viab-real}

\pt{%
Como responsável pela área de dados, o autor explorou os dados da operação no ambiente da empresa, apenas com consultas agregadas, sem dados pessoais e com as unidades anonimizadas. O painel cobriu doze meses comuns às três fontes principais, com algumas centenas de acidentes relatados ligados a veículos com telemetria e cerca de duas mil combinações de veículo e mês. As hipóteses foram testadas com os procedimentos da Seção~\ref{sec:testes} numa versão simplificada: exposição pelo tempo de viagem, e modelo de H2 com tipo de uso e efeito aleatório de veículo, mas sem efeito de unidade nem de calendário. Os números vêm de \texttt{analise/real/testes.py} e \texttt{r/h2\_misto.R}; a Figura~\ref{fig:real} e o texto mostram só valores relativos e arredondados.
}
\en{%
As head of the data area, the author explored the operational data inside the company's environment, using only aggregate queries, with no personal data and with sites anonymized. The panel covered twelve months common to the three main sources, with a few hundred reported accidents linked to vehicles with telematics and about two thousand vehicle-month combinations. The hypotheses were tested with the procedures of Section~\ref{sec:testes} in a simplified version: exposure by trip time, and an H2 model with use type and a vehicle random effect, but no site or calendar effect. Numbers come from \texttt{analise/real/testes.py} and \texttt{r/h2\_misto.R}; Figure~\ref{fig:real} and the text show relative and rounded values only.
}

\begin{figure}[!htbp]
\centering
\includegraphics[width=0.9\textwidth]{\arqfig{fig_real}}
\caption{\tr%
  {Exploração dos dados da operação. (a)~Taxa de acidentes de cada tipo de uso relativa à do caminhão de manobra (\textit{spotter}), por milha e por hora de viagem. (b)~Gráfico de funil: taxa de cada unidade relativa à média, bruta e por Bayes empírico, contra a exposição relativa à mediana; linhas: limites de 95\% e 99,8\%. (c)~Razão de taxas de acidente por desvio padrão da taxa de cada tipo de evento de condução, sem ajuste e com tipo de uso e efeito aleatório de veículo; \corpb{faixa azul-clara}{faixa cinza}: margem de equivalência. (d)~Proporção de acidentes relatados com detecção automática de colisão no mesmo veículo em até dois dias, por tipo de uso, e proporção de detecções descartadas na revisão. Barras de erro: intervalos de 95\%.}%
  {Exploration of operational data. (a)~Accident rate of each use type relative to the yard truck (spotter), per mile and per trip hour. (b)~Funnel plot: each site's rate relative to the mean, raw and empirical Bayes, against exposure relative to the median; lines: 95\% and 99.8\% limits. (c)~Accident rate ratio per standard deviation of each driving event rate, unadjusted and with use type and vehicle random effect; \corpb{light blue band}{grey band}: equivalence margin. (d)~Share of reported accidents with an automatic crash detection on the same vehicle within two days, by use type, and share of detections dismissed in review. Error bars: 95\% intervals.}}
\label{fig:real}
\end{figure}

\pt{%
\textbf{Exposição e tipo de uso.} O caminhão de manobra anda pouco e manobra muito: por milha, sua taxa de acidentes é mais de vinte vezes a dos veículos de uso local; por hora, cerca de três vezes (Figura~\ref{fig:real}a). A escolha da medida de exposição muda a ordem de grandeza da comparação. \textbf{Implicação:} o tempo de viagem é a exposição principal, e o tipo de uso entra como ajuste obrigatório.
}
\en{%
\textbf{Exposure and use type.} The yard truck travels little and maneuvers a lot: per mile, its accident rate is more than twenty times that of local-use vehicles; per hour, about three times (Figure~\ref{fig:real}a). The choice of exposure measure changes the order of magnitude of the comparison. \textbf{Implication:} trip time is the main exposure, and use type enters as a mandatory adjustment.
}

\pt{%
\textbf{H1.} As taxas variaram entre unidades além do acaso: a homogeneidade foi rejeitada ($p < 0{,}001$), e cerca de um sexto das unidades ficou fora do limite de 95\% do funil, contra 5\% esperados (H1b; Figura~\ref{fig:real}b). A taxa por Bayes empírico previu a taxa do semestre seguinte com erro menor que a taxa bruta ($p \approx 0{,}02$; H1a). A contagem bruta acompanhou o tamanho da unidade (correlação de postos de cerca de 0,8), e a taxa por Bayes empírico não (intervalo de 90\% dentro de $\pm 0{,}3$; H1c). Ainda assim, a ordenação das unidades mudou muito de um semestre para o outro. \textbf{Implicação:} os resultados da simulação se confirmam; a comunicação deve ser por intervalos e grupos.
}
\en{%
\textbf{H1.} Rates varied across sites beyond chance: homogeneity was rejected ($p < 0.001$), and about one sixth of the sites fell outside the funnel's 95\% limit, against 5\% expected (H1b; Figure~\ref{fig:real}b). The empirical Bayes rate predicted the next semester's rate with smaller error than the raw rate ($p \approx 0.02$; H1a). The raw count tracked site size (rank correlation of about 0.8), whereas the empirical Bayes rate did not (90\% interval within $\pm 0.3$; H1c). Even so, the ordering of sites changed considerably from one semester to the next. \textbf{Implication:} the simulation results are confirmed; communication should be through intervals and groups.
}

\pt{%
\textbf{H2.} Sem ajuste, alguns tipos de evento pareciam associados a acidentes. Com tipo de uso e efeito de veículo, nenhum dos dez tipos rejeitou $H_0$ após a correção de Holm, e o modelo com telemetria teve escore fora da amostra ligeiramente pior que o modelo sem ela (H2a). Sete tipos tiveram efeito declarado desprezível; frenagem brusca, excesso grave de velocidade e parada incompleta ficaram inconclusivos (H2b; Figura~\ref{fig:real}c). \textbf{Implicação:} a associação aparente vinha do tipo de operação; o modelo completo e mais meses de dados vão decidir os três casos inconclusivos.
}
\en{%
\textbf{H2.} Without adjustment, some event types seemed associated with accidents. With use type and vehicle effect, none of the ten types rejected $H_0$ after Holm correction, and the model with telematics had a slightly worse out-of-sample score than the model without it (H2a). Seven types had their effect declared negligible; harsh braking, severe speeding and rolling stops were inconclusive (H2b; Figure~\ref{fig:real}c). \textbf{Implication:} the apparent association came from the type of operation; the full model and more months of data will settle the three inconclusive cases.
}

\pt{%
\textbf{H3 e detecção automática.} Menos de um décimo dos acidentes relatados teve alguma detecção automática no mesmo veículo em até dois dias, proporção igual nos três tipos de uso, e mais de quatro quintos das detecções foram descartadas como falso alarme na revisão (Figura~\ref{fig:real}d). Os sensores foram calibrados para impactos de estrada, e o acidente típico do pátio é de baixa velocidade. \textbf{Implicação:} a detecção automática não serve como segunda fonte no pátio; a captura e recaptura passa a usar a manutenção (Seção~\ref{sec:met-subnotificacao}).
}
\en{%
\textbf{H3 and automatic detection.} Fewer than one tenth of reported accidents had any automatic detection on the same vehicle within two days, an equal proportion across the three use types, and more than four fifths of detections were dismissed as false alarms in review (Figure~\ref{fig:real}d). The sensors are calibrated for road impacts, and the typical yard accident is low-speed. \textbf{Implication:} automatic detection does not work as a second source in the yard; capture-recapture now uses maintenance (Section~\ref{sec:met-subnotificacao}).
}

\pt{%
\textbf{Modelo e extensão.} O efeito aleatório de veículo foi significativo, e os resíduos quantílicos simulados não indicaram superdispersão residual nem excesso de zeros, o que sustenta a binomial negativa com efeito de veículo, sem inflação de zeros, como ponto de partida. As horas semanais trabalhadas foram previstas com erro de cerca de 6\% num horizonte de oito semanas.
}
\en{%
\textbf{Model and extension.} The vehicle random effect was significant, and simulated quantile residuals showed no residual overdispersion or excess zeros, which supports the negative binomial with a vehicle effect and no zero inflation as the starting point. Weekly hours worked were forecast with an error of about 6\% over an eight-week horizon.
}

\pt{%
São resultados exploratórios, de um único ano e com uma versão simplificada do modelo; servem para orientar o desenho, não como conclusão do trabalho.
}
\en{%
These are exploratory results, from a single year and with a simplified version of the model; they guide the design and are not conclusions of the work.
}

\section{\tr{Síntese}{Summary}}

\pt{%
\begin{enumerate}
  \item O projeto é viável na escala esperada se o desfecho incluir todas as colisões relatadas (centenas no período).
  \item O efeito aleatório de caminhão não é opcional: sem ele, a inferência sobre telemetria fica otimista. Nos dados reais, sua variância foi significativa.
  \item O tipo de uso e a exposição por tempo são obrigatórios: sem eles, a comparação entre veículos e a associação com a telemetria refletem o tipo de operação.
  \item Ranking de unidades é ruidoso: comunicar com intervalos e grupos, e destacar a troca de contagem por taxa ajustada. Os dados reais confirmaram a variação real entre unidades e a vantagem da taxa com encolhimento.
  \item A subnotificação por unidade só é estimável com encolhimento, e a segunda fonte tem de ser a manutenção, não a detecção automática.
  \item A extensão (f) foi reposicionada para o planejamento de escala.
\end{enumerate}
}
\en{%
\begin{enumerate}
  \item The project is feasible at the expected scale if the outcome includes all reported collisions (hundreds in the period).
  \item The truck random effect is not optional: without it, inference on telematics is optimistic. In the real data, its variance was significant.
  \item Use type and time-based exposure are mandatory: without them, the comparison between vehicles and the association with telematics reflect the type of operation.
  \item Site ranking is noisy: communicate with intervals and groups, and stress replacing counts with adjusted rates. The real data confirmed real variation across sites and the advantage of the shrinkage rate.
  \item Site-level underreporting is only estimable with shrinkage, and the second source must be maintenance, not automatic detection.
  \item Extension (f) was repositioned toward staffing planning.
\end{enumerate}
}

\chapter{\tr{Análise preditiva}{Predictive analysis}}
\label{cap:preditiva}
\label{cap:resultados}

\pt{%
Este capítulo ajusta e avalia os modelos de risco e testa as hipóteses, partindo das distribuições descritas no Capítulo~\ref{cap:descritiva}. Depende dos dados completos da operação, depois da autorização formal da empresa e da retomada da sincronização da telemetria. A estrutura abaixo segue as etapas do Capítulo~\ref{cap:metodos}, com o que cada seção vai mostrar. Todos os números serão apresentados como valores relativos, porcentagens arredondadas e ordens de grandeza.
}
\en{%
This chapter fits and evaluates the risk models and tests the hypotheses, starting from the distributions described in Chapter~\ref{cap:descritiva}. It depends on the full operational data, after the company's formal authorization and the resumption of the telematics sync. The structure below follows the stages of Chapter~\ref{cap:metodos}, with what each section will show. All numbers will be presented as relative values, rounded percentages and orders of magnitude.
}

\section{\tr{Painel final e qualidade da ligação}{Final panel and linkage quality}}
\label{sec:res-painel}

\pt{%
Atualização das métricas do Capítulo~\ref{cap:descritiva} com o período completo; qualidade da ligação (precisão e cobertura na amostra revisada) e efeito de cada regra de limpeza no número de relatos. \preencher{Rodar de novo \texttt{analise/real/descritiva.py} e \texttt{r/descritiva.R} com o período completo.}
}
\en{%
Update of the Chapter~\ref{cap:descritiva} metrics with the full period; linkage quality (precision and coverage in the reviewed sample) and the effect of each cleaning rule on the number of reports. \preencher{Rerun \texttt{analise/real/descritiva.py} and \texttt{r/descritiva.R} with the full period.}
}

\section{\tr{Seleção de modelo (objetivo b)}{Model selection (objective b)}}
\label{sec:res-modelos}

\pt{%
Comparação dos modelos M0 a M6 (Tabela~\ref{tab:modelos}) por AIC e por escore logarítmico na validação temporal e por unidade; testes das variâncias dos efeitos aleatórios; verificação do \textit{offset} ($\hat{\theta}$ perto de 1?); diagnóstico por resíduos simulados do modelo escolhido. \preencher{Tabela de comparação e figura dos resíduos.}
}
\en{%
Comparison of models M0 to M6 (Table~\ref{tab:modelos}) by AIC and by log score in temporal and site validation; tests of random effect variances; offset check ($\hat{\theta}$ near 1?); simulated residual diagnostics of the chosen model. \preencher{Comparison table and residual figure.}
}

\section{\tr{Comparação entre unidades (H1, objetivo e)}{Site comparison (H1, objective e)}}
\label{sec:res-h1}

\pt{%
H1a, H1b e H1c com os procedimentos da Seção~\ref{sec:testes}; gráfico de funil com o modelo completo; grupos de risco e intervalos de posição; estabilidade da classificação entre semestres. \preencher{Tabela de testes de H1, gráfico de funil e figura de intervalos de posição.}
}
\en{%
H1a, H1b and H1c with the procedures of Section~\ref{sec:testes}; funnel plot with the full model; risk groups and position intervals; stability of the classification across semesters. \preencher{H1 test table, funnel plot and position interval figure.}
}

\section{\tr{Telemetria (H2, objetivo c)}{Telematics (H2, objective c)}}
\label{sec:res-h2}

\pt{%
Razões de taxas por tipo de evento, sem ajuste e no modelo completo, com classificação em associado, desprezível ou inconclusivo; diferença de escore fora da amostra entre M4 e M3; concentração de acidentes nos veículos de maior risco previsto, com e sem telemetria. \preencher{Tabela de H2 e figura das razões de taxas com a margem de equivalência.}
}
\en{%
Rate ratios by event type, unadjusted and in the full model, classified as associated, negligible or inconclusive; out-of-sample score difference between M4 and M3; accident concentration in the vehicles with the highest predicted risk, with and without telematics. \preencher{H2 table and figure of rate ratios with the equivalence margin.}
}

\section{\tr{Subnotificação (H3, objetivo d)}{Underreporting (H3, objective d)}}
\label{sec:res-h3}

\pt{%
Cobertura da detecção automática de colisão no período completo; pareamento entre ordens de dano e relatos (distribuição dos intervalos, janela escolhida); probabilidade de relato geral e por unidade, com encolhimento; testes de H3a e H3b; total de acidentes estimado por captura e recaptura, por estrato, como limite. \preencher{Tabela de H3 e figura das probabilidades de relato por unidade.}
}
\en{%
Coverage of automatic crash detection over the full period; matching between damage orders and reports (interval distribution, chosen window); overall and site reporting probability, with shrinkage; tests of H3a and H3b; total accidents estimated by capture-recapture, by stratum, as a bound. \preencher{H3 table and figure of site reporting probabilities.}
}

\section{\tr{Análises de sensibilidade}{Sensitivity analyses}}
\label{sec:res-sensibilidade}

\pt{%
As conclusões de H1 a H3 sob cada variação da Seção~\ref{sec:met-sensibilidade}. \preencher{Tabela-resumo: conclusão principal e se ela se mantém em cada variação.}
}
\en{%
The conclusions of H1 to H3 under each variation in Section~\ref{sec:met-sensibilidade}. \preencher{Summary table: main conclusion and whether it holds under each variation.}
}

\section{\tr{Previsão da exposição (objetivo f, opcional)}{Exposure forecasting (objective f, optional)}}
\label{sec:res-previsao}

\pt{%
Desempenho dos três métodos por horizonte (CRPS, perda quantílica, cobertura); projeção de acidentes por região e trimestre e fração da variância devida à exposição. \preencher{Só se a extensão for feita.}
}
\en{%
Performance of the three methods by horizon (CRPS, quantile loss, coverage); accident projection by region and quarter and share of variance due to exposure. \preencher{Only if the extension is carried out.}
}

\chapter{\tr{Análise prescritiva}{Prescriptive analysis}}
\label{cap:prescritiva}

\pt{%
\preencher{Capítulo planejado; depende dos resultados do Capítulo~\ref{cap:preditiva}.} A análise preditiva estima quanto risco há e onde. A prescritiva transforma essas estimativas em regras de decisão para a gestão de segurança e de operação, e avalia as regras pelo que teriam produzido no passado. São três decisões, em ordem de prioridade.
}
\en{%
\preencher{Planned chapter; depends on the results of Chapter~\ref{cap:preditiva}.} The predictive analysis estimates how much risk there is and where. The prescriptive analysis turns these estimates into decision rules for safety and operations management, and evaluates the rules by what they would have produced in the past. There are three decisions, in order of priority.
}

\section{\tr{Onde agir: priorização de unidades}{Where to act: site prioritization}}
\label{sec:presc-unidades}

\pt{%
A gestão escolhe periodicamente um pequeno número de unidades para treinamento, inspeção ou supervisão reforçada. Hoje a escolha segue a contagem bruta, que o Capítulo~\ref{cap:descritiva} mostrou medir sobretudo o tamanho. Três regras serão comparadas, sempre escolhendo o mesmo número $k$ de unidades: (i) maior contagem bruta; (ii) maior taxa ajustada $\exp(\hat{u}_s)$; (iii) maior excesso esperado de acidentes, $E_s(\hat{\lambda}_s - \lambda_{\text{ref}})$, que combina risco e tamanho e responde à pergunta ``onde uma redução de risco evitaria mais acidentes''. Variantes de (ii) e (iii) usam a probabilidade a posteriori de a unidade estar acima da referência, para não premiar estimativas instáveis \parencite{hauer1997,goldstein1996}.
}
\en{%
Management periodically chooses a small number of sites for training, inspection or reinforced supervision. Today the choice follows the raw count, which Chapter~\ref{cap:descritiva} showed to mostly measure size. Three rules will be compared, always choosing the same number $k$ of sites: (i) highest raw count; (ii) highest adjusted rate $\exp(\hat{u}_s)$; (iii) highest expected excess accidents, $E_s(\hat{\lambda}_s - \lambda_{\text{ref}})$, which combines risk and size and answers the question ``where would a risk reduction prevent the most accidents''. Variants of (ii) and (iii) use the posterior probability that the site is above the reference, so as not to reward unstable estimates \parencite{hauer1997,goldstein1996}.
}

\pt{%
A avaliação é retrospectiva: escolher as unidades com os dados de um semestre e medir, no semestre seguinte, quantos acidentes em excesso ao esperado ocorreram nas escolhidas. Como não há intervenção aleatorizada, a avaliação mede a capacidade de cada regra de apontar o risco futuro, não o efeito de agir sobre ele. \preencher{Resultados.}
}
\en{%
Evaluation is retrospective: choose sites with one semester's data and measure, in the next semester, how many accidents in excess of expected occurred in the chosen ones. Since there is no randomized intervention, evaluation measures each rule's ability to point to future risk, not the effect of acting on it. \preencher{Results.}
}

\section{\tr{Quanto alocar: escala de motoristas por turno}{How much to allocate: driver staffing by shift}}
\label{sec:presc-escala}

\pt{%
É a extensão (f). A previsão probabilística da demanda de movimentos por unidade e turno vira um número de motoristas. Em vez de somar uma folga fixa à previsão média, o número é escolhido pelo quantil da previsão que equilibra o custo de faltar e o de sobrar motorista, como no problema do jornaleiro:
}
\en{%
This is extension (f). The probabilistic forecast of move demand per site and shift becomes a number of drivers. Instead of adding a fixed buffer to the mean forecast, the number is chosen by the forecast quantile that balances the cost of having too few and too many drivers, as in the newsvendor problem:
}
\begin{equation}
\label{eq:jornaleiro}
n^{*} = \left\lceil \frac{F^{-1}(q^{*})}{c} \right\rceil, \qquad q^{*} = \frac{c_{\text{falta}}}{c_{\text{falta}} + c_{\text{sobra}}},
\end{equation}
\pt{%
em que $F$ é a distribuição prevista de movimentos no turno, $c$ é a capacidade de um motorista no turno (movimentos por motorista-hora vezes horas efetivas) e $c_{\text{falta}}$ e $c_{\text{sobra}}$ são os custos unitários de faltar e de sobrar um motorista. A capacidade deve ser estimada na mesma unidade da comparação com a escala atual (motorista-horas, não pessoas distintas). As regras serão avaliadas por \textit{backtest} de decisão: escalar com a previsão feita no passado e contar os turnos que teriam ficado curtos ou com sobra, e o custo de cada um. Se a carga de trabalho por turno se mostrar associada ao risco, o custo de faltar motorista passa a incluir o risco de acidente. \preencher{Depende da escala da folha por unidade e turno; ver \texttt{notas/10}.}
}
\en{%
where $F$ is the forecast distribution of moves in the shift, $c$ is one driver's capacity in the shift (moves per driver-hour times effective hours) and $c_{\text{falta}}$ and $c_{\text{sobra}}$ are the unit costs of being one driver short and over. Capacity must be estimated in the same unit as the comparison with current staffing (driver-hours, not distinct people). Rules will be evaluated by a decision backtest: staff with the forecast made in the past and count the shifts that would have been short or over, and the cost of each. If shift workload turns out to be associated with risk, the cost of being short includes accident risk. \preencher{Depends on payroll scheduling by site and shift.}
}

\section{\tr{Quando inspecionar: alocação de auditorias}{When to inspect: audit allocation}}
\label{sec:presc-auditorias}

\pt{%
As auditorias de segurança por unidade (Capítulo~\ref{cap:descritiva}) consomem tempo de supervisores. Uma regra simples prioriza as unidades com maior risco estimado e maior incerteza sobre ele, onde uma auditoria tem mais a informar, e a compara com o calendário fixo atual. \preencher{Opcional; avaliar se a nota da auditoria varia o bastante para ser útil.}
}
\en{%
Site safety audits (Chapter~\ref{cap:descritiva}) take up supervisors' time. A simple rule prioritizes the sites with the highest estimated risk and the greatest uncertainty about it, where an audit has the most to reveal, and compares it with the current fixed schedule. \preencher{Optional; assess whether the audit score varies enough to be useful.}
}

\chapter{\tr{Conclusão}{Conclusion}}
\label{cap:conclusao}

\pt{%
\preencher{Versão provisória, escrita antes dos resultados finais. Revisar depois do Capítulo~\ref{cap:resultados}.}
}
\en{%
\preencher{Provisional version, written before the final results. Revise after Chapter~\ref{cap:resultados}.}
}

\section{\tr{Síntese do trabalho}{Summary of the work}}

\pt{%
O trabalho propôs medir o risco de acidentes em operações de pátio de forma comparável entre veículos e unidades, ajustando pela exposição e pelo tipo de operação, com a incerteza explícita, e testar duas promessas da telemetria de estrada nesse contexto: os eventos de condução como indicador de risco e a detecção automática de colisão como fonte de acidentes. Os estudos preliminares já mostraram que o desenho é viável na escala da operação, desde que o desfecho inclua todas as colisões relatadas e o modelo tenha efeito de veículo, e que a escolha da exposição e o ajuste pelo tipo de uso mudam as conclusões. \preencher{Resumir os resultados finais de H1, H2 e H3.}
}
\en{%
The work set out to measure accident risk in yard operations comparably across vehicles and sites, adjusting for exposure and type of operation, with explicit uncertainty, and to test two promises of road telematics in this setting: driving events as a risk indicator and automatic crash detection as a source of accidents. The preliminary studies already showed that the design is feasible at the operation's scale, provided the outcome includes all reported collisions and the model has a vehicle effect, and that the choice of exposure and the adjustment for use type change the conclusions. \preencher{Summarize the final results for H1, H2 and H3.}
}

\section{\tr{Contribuições}{Contributions}}

\pt{%
\begin{itemize}
  \item \textbf{Para a gestão de segurança:} um indicador de risco por unidade que não se confunde com o tamanho da unidade e que prevê o período seguinte melhor que a taxa bruta, comunicado por grupos de risco e gráficos de funil em vez de rankings.
  \item \textbf{Para a literatura de segurança:} evidência, num tipo de operação sem estudo próprio, sobre o limite da telemetria de estrada quando aplicada ao pátio; um resultado nulo bem delimitado, por testes de equivalência, é informativo, porque indica quais indicadores não vale a pena acompanhar.
  \item \textbf{Para a prática estatística:} um estudo de simulação calibrado em dados públicos que mostra, para frotas de algumas centenas de veículos, o número de acidentes necessário e o custo de omitir o efeito de veículo.
\end{itemize}
\preencher{Ajustar às conclusões finais.}
}
\en{%
\begin{itemize}
  \item \textbf{For safety management:} a site risk indicator that is not confounded with site size and predicts the next period better than the raw rate, communicated through risk groups and funnel plots instead of rankings.
  \item \textbf{For the safety literature:} evidence, in a type of operation with no study of its own, on the limits of road telematics when applied to the yard; a well-bounded null result, through equivalence tests, is informative, because it shows which indicators are not worth tracking.
  \item \textbf{For statistical practice:} a simulation study calibrated on public data showing, for fleets of a few hundred vehicles, the number of accidents required and the cost of omitting the vehicle effect.
\end{itemize}
\preencher{Adjust to the final conclusions.}
}

\section{\tr{Limitações}{Limitations}}

\pt{%
\begin{itemize}
  \item \textbf{Estudo observacional.} As associações com a telemetria não são efeitos causais: motorista, unidade e tipo de manobra podem gerar ao mesmo tempo eventos e acidentes. O trabalho mede valor preditivo, não o efeito de mudar o comportamento de condução.
  \item \textbf{Uma empresa e um período curto.} Os dados vêm de uma única empresa, com pouco mais de um ano comum às fontes; a generalização para outras operações de pátio depende de processos de relato e de frotas semelhantes.
  \item \textbf{Desfecho sujeito a relato.} O desfecho principal é o acidente relatado. A Etapa~4 estima a subnotificação, mas só entre os acidentes que danificam o próprio caminhão, e o total estimado é um limite sob hipóteses sobre a dependência entre as fontes.
  \item \textbf{Ligação imperfeita.} Parte dos relatos tem veículo e unidade digitados em texto livre; erros de ligação afetam tanto o desfecho por veículo quanto o pareamento da captura e recaptura.
  \item \textbf{Telemetria de estrada.} Os eventos de condução usam limiares calibrados para a estrada; um resultado nulo diz respeito a esses eventos, não a qualquer medida de condução no pátio.
\end{itemize}
}
\en{%
\begin{itemize}
  \item \textbf{Observational study.} Associations with telematics are not causal effects: driver, site and type of maneuver can generate both events and accidents. The work measures predictive value, not the effect of changing driving behavior.
  \item \textbf{One company and a short period.} The data come from a single company, with a little over one year common to the sources; generalizing to other yard operations depends on similar reporting processes and fleets.
  \item \textbf{Outcome subject to reporting.} The main outcome is the reported accident. Stage~4 estimates underreporting, but only among accidents that damage the truck itself, and the estimated total is a bound under assumptions about dependence between sources.
  \item \textbf{Imperfect linkage.} Some reports have vehicle and site typed as free text; linkage errors affect both the per-vehicle outcome and the capture-recapture matching.
  \item \textbf{Road telematics.} Driving events use thresholds calibrated for the road; a null result concerns these events, not any measure of yard driving.
\end{itemize}
}

\section{\tr{Trabalhos futuros}{Future work}}

\pt{%
\begin{itemize}
  \item \textbf{Exposição própria do pátio.} Usar o número de movimentos de carreta registrados pelo sistema de gestão de pátio como exposição, se ele puder ser ligado ao caminhão; é uma medida mais próxima da manobra que o tempo de viagem.
  \item \textbf{Eventos próprios do pátio.} Definir, com os dados brutos dos sensores, eventos de condução com limiares de baixa velocidade (ré, engate, aproximação da doca) e testar se eles têm o sinal que os eventos de estrada não têm.
  \item \textbf{Mais fontes para a subnotificação.} Acrescentar uma terceira fonte (por exemplo, reclamações de clientes por dano à carreta), o que permite modelar a dependência entre fontes com modelos log-lineares \parencite{iwgdmf1995}.
  \item \textbf{Intervenções.} Usar o indicador ajustado para avaliar treinamentos e mudanças de supervisão, com o Bayes empírico para corrigir a regressão à média na escolha das unidades \parencite{hauer1997}.
\end{itemize}
}
\en{%
\begin{itemize}
  \item \textbf{Yard-specific exposure.} Use the number of trailer moves recorded by the yard management system as exposure, if it can be linked to the truck; it is a measure closer to maneuvering than trip time.
  \item \textbf{Yard-specific events.} Define, from raw sensor data, driving events with low-speed thresholds (reversing, coupling, dock approach) and test whether they carry the signal that road events do not.
  \item \textbf{More sources for underreporting.} Add a third source (for example, client claims for trailer damage), which allows modeling dependence between sources with log-linear models \parencite{iwgdmf1995}.
  \item \textbf{Interventions.} Use the adjusted indicator to evaluate training and supervision changes, with empirical Bayes to correct regression to the mean in site selection \parencite{hauer1997}.
\end{itemize}
}


\begingroup
\singlespacing
\setlength{\bibitemsep}{\baselineskip}
\printbibliography[title={\tr{Referências}{References}}, heading=bibintoc]
\endgroup

\end{document}
